% concatenated from STRONG_STABLE_SETS.tex
\documentclass[a4paper,english,reqno]{amsart}

\usepackage{amssymb,amsthm, amsmath,amsfonts}
\usepackage{bbm}
\usepackage{graphicx}
\graphicspath{{figs/}}
\usepackage{subfigure}
\usepackage{microtype}
\usepackage{booktabs}
\usepackage{enumerate}
\usepackage{pgf}
\usepackage[latin1]{inputenc}
\usepackage{verbatim}
\usepackage{cite}
\usepackage{url}
\usepackage{xcolor}

\usepackage[cmtip,all]{xy}
\xyoption{arc}

\usepackage{pgf,tikz}
\usetikzlibrary{arrows,shapes,calc,positioning}
\usepackage{manfnt}
\usepackage{tikz-cd}
\usepackage{url}
\usepackage{microtype}

\usepackage[font=footnotesize]{caption}

\usepackage{xcolor}
\usepackage{hyperref}
%\usepackage{enumitem}

%some operators
%\DeclareMathOperator{\RR}{\mathbb{R}}



\newcommand{\Z}{\mathbb{Z}}
\newcommand{\R}{\mathbb{R}}
\newcommand{\C}{\mathbb{C}}
\renewcommand{\H}{\mathbb{H}}
\newcommand{\D}{\mathbb{D}}
\newcommand{\G}{\mathbb{G}}
\newcommand{\M}{\mathbb{M}}
\newcommand{\N}{\mathbb{N}} %este lo añadi yo
\newcommand{\F}{\mathcal{F}}
\newcommand{\scirc}{\raisebox{0.2ex}{$\scriptscriptstyle \circ$}}
\newcommand{\bs}{\backslash}
\newcommand{\vf}[1]{\frac{\partial~}{\partial #1}}
\newcommand{\secf}[1]{\frac{\partial^2~}{\partial #1^2}}

\newcommand{\blue}{\color{blue}}
\newcommand{\red}{\color{red}}
\newcommand{\green}{\color{green}}


%     Update the information and uncomment if AMS is not the copyright
%     holder.
%\copyrightinfo{2009}{American Mathematical Society}

\newtheorem{theorem}{Theorem}%[section]
\newtheorem{lemma}{Lemma}
\newtheorem{proposition}{Proposition}
\newtheorem{corollary}{Corollary}
\newtheorem{property}{Property} %este lo anado yo
\newtheorem*{theoremA*}{Main Theorem}
\newtheorem*{theoremB*}{Main Theorem for a Flute}


\theoremstyle{definition}
\newtheorem{definition}[theorem]{Definition}
\newtheorem{example}[theorem]{Example}
\newtheorem{notation}[theorem]{Notation} %este lo anado yo
%\newtheorem{xca}[theorem]{Exercise}
\newtheorem{remark}{Remark}



\numberwithin{equation}{section}

\begin{document}
	
	\title[Bellis strong stable sets on 
	infinite hyperbolic surfaces]%
	{Bellis strong stable sets on 
		infinite hyperbolic surfaces}
	
	%    Only \author and \address are required; other information is
	%    optional.  Remove any unused author tags.
	
	%    author one information
	% \author[short version for running head]{name for top of paper}
	%\author{}
	%\address{}
	%\curraddr{}
	%\email{}
	%%\thanks{}
	%\author[F. Alcalde]{Fernando Alcalde Cuesta}
	%\address{CITMaga, Universidad de Santiago de Compostela, E-15782, Santiago de Compostela, Spain.
		%}
	%\email{fernando.alcalde@usc.es}
	%\thanks{}
	%
	%%2th author
	%\author[\'A. Carballido]{\'Alvaro Carballido Costas}
	%\address{CITMaga, Universidad de Santiago de Compostela, Santiago de Compostela, Spain.}
	%\email{alvaro.carballido.costas@usc.es}
	
	%3th author

	%\author[A.Bellis]{Alexandre Bellis}
	%\author[f. Dal'Bo]{Fran\c coise Dal'Bo}
	
	\author[S. Burniol, F. Dal'Bo, S. Herrero]{ Sergi Burniol Clotet, Fran�oise Dal'Bo, Sergio Herrero Vila}	
	\address{IMERL, CMAT, Universidad de la Rep�blica, IRL-CNRS-IFUMI 2030, Uruguay}
	\email{sergi.burniol@gmail.com}
	\address{IRMAR, Universit\'e de Rennes, France, IRL-CNRS-IFUMI 2030, Uruguay}
	\email{francoise.dalbo@univ-rennes.fr}
	\address{IRMAR, Universit\'e de Rennes, Rennes, France}
	\email{sergioherrerovila@outlook.com}
	%    \subjclass is required.
	\subjclass[2020]{Primary 37D40; Secondary 53D25, 53C22.}
	
	\date{01/04/2026}
	
	
	\dedicatory{}
	
	\begin{abstract}		
		We provide a corrected proof of a theorem of A. Bellis on strong stable sets in the unit tangent bundle of certain hyperbolic surfaces. The theorem states that, for vectors whose geodesic rays encounter arbitrarily short closed geodesics, the strong stable set in the dynamical sense does not coincide with the associated horocyclic orbit. The proof is based on Bellis' idea of constructing geodesic rays that wind around infinitely many closed geodesics.
	\end{abstract}
	
	\maketitle
	
	%	\textcolor{blue}{INTRODUCIR UNA PEQUE�A INTRODUCCION DE LO QUE HACEMOS, AQUI ON EN EL ABSTRACT}.
	%	
	\section{Introduction}
	The purpose of this paper is to give a corrected proof of a theorem stated in the PhD thesis of A. Bellis \cite[Th�or�me E]{Bellis} on strong stable sets of the geodesic flow $g_t$ on the unit tangent bundle $T^1 S$ of a hyperbolic surface $S$. We consider the Sasaki distance $d_{Sa}$ on $T^1 S$ induced by the hyperbolic metric of $S$. The strong stable set $W^{ss}u$ of a vector $u\in T^1S$ is defined as 
	\[
	W^{ss}u = \{ v\in T^1S \, | \, \lim\limits_{t\to +\infty} d_{Sa}(g_tu,g_tv)=0  \}.
	\]
	Let $h_t$ denote the horocyclic flow on $T^1 S$.
	
	The horocyclic orbit $h_{\mathbb{R}} u$ of any vector $u\in T^1S$ is always contained in $W^{ss} u$ \cite{Dalbo}. Bellis proves that if the infimum of the injectivity radius along the geodesic ray $u(\mathbb{R}^+)$ is positive, then $h_{\mathbb{R}} u = W^{ss} u$ \cite[Th�or�me D]{Bellis}. When $S$ contains arbitrarily small closed geodesics --in particular, $S$ is of infinite type--, Bellis detects vectors $u$ for which $h_{\mathbb{R}} u$ and  $W^{ss} u$ are different.
	
	\begin{theoremA*}\label{thE} \cite[Th�or�me E]{Bellis}
		Let $S$ be a hyperbolic surface and $u\in T^1S$. If the geodesic ray $u(\mathbb{R}^+)$ meets an infinite sequence of closed geodesics whose lengths tend to zero, then
		\begin{enumerate}
			\item $h_{\mathbb{R}} u\subsetneq W^{ss}u$,
			\item the set $W^{ss}u $ is an uncountable union of horocyclic orbits $h_{\mathbb{R}} v_i$.
		\end{enumerate} 
	\end{theoremA*}
	
	The proof of Bellis contains several errors. In particular, there is an important gap in the proof of Lemma 4.3.1 in \cite{Bellis}. We are only able to prove a weaker statement (Definition \ref{nested_sequence}, Proposition \ref{prop3}). We present a new argument to fill the gap in the proof (Section \ref{end_proof}). When $S$ is a fine flute surface (see \cite[Section 1.4.1]{Bellis}), our argument is not needed, and Bellis' proof can be considered essentially correct.
	
	We also have simplified and corrected the proofs of Proposition 4.1.5 and Lemma 4.3.4 in \cite{Bellis}. We introduce a distance $d_1$ on $T^1S$ involving only the distance on $S$ and avoiding parallel transport, which appears to be a source of errors in the original arguments. We prove in the appendix that $d_1$ and $d_{Sa}$ define the same strong stable sets.
	
	The paper is self-contained. In Section 3, we formalize the notion of winding around a closed geodesic, which is the key tool used by Bellis. Section 4 is devoted to the proof of the Main Theorem.
	
	%    "Communicated by" -- provide editor's name; required.
	%\commby{}
	
	%    Abstract is required.
	
	\maketitle
		
%	\textcolor{blue}{INTRODUCIR UNA PEQUE�A INTRODUCCION DE LO QUE HACEMOS, AQUI ON EN EL ABSTRACT}.
%	
	\section*{}
	\textbf{Acknowledgment :} The third author wishes to express his sincere gratitude to the Department of Mathematics in Montevideo for its warm hospitality during September 2023.
		
		
	
	
	
	
%	\begin{theoremB*}\label{toytheorem}
%		Let $S_F$ be a hyperbolic flute (see definition \ref{def_flute}).
%		Choose $\tilde u_0\in T^1 \H^2$, with $\tilde u_0(+\infty)=\infty$ and denote by $u_0$ its projection onto $T^1 S_F$. If the sequence $(\ell(\gamma_n))_{n\in \N}$ goes to zero, then
%		\begin{enumerate}
%			\item $h_\R u\subsetneq W^{ss}u_0.$,
%			\item the set $W^{ss}u_0 $ is an uncountable union of horocyclic trajectories.
%		\end{enumerate} 
%	\end{theoremB*}

\section{Preliminaries}

We denote by $\mathbb H^2=\{z\in\mathbb{C}:\operatorname{Im} (z)>0\}$ the hyperbolic plane endowed with the hyperbolic distance $d$, and we will denote by $
\partial \mathbb H^2=\mathbb R\cup\{\infty\}
$
 its boundary at infinity.


Recall that the orientation-preserving isometries of $\mathbb H^2$ are precisely the
M�bius transformations
\[
z\longmapsto \frac{az+b}{cz+d},
\qquad
\begin{pmatrix}
	a&b\\ c&d
\end{pmatrix}\in \operatorname{PSL}(2,\mathbb R).
\]
%Hence, $\operatorname{PSL}(2,\mathbb R)$ identifies with the group of orientation-preserving
%isometries of $\mathbb H^2$.

An isometry $\gamma\in \operatorname{PSL}(2,\mathbb R)$ is called \emph{hyperbolic} if it has two fixed
points on $\partial\mathbb H^2$, denoted by $\gamma^+$ and $\gamma^-$, called respectively
the attractive and repulsive fixed points of $\gamma$. Its translation length along the geodesic $(\gamma^-,\gamma^+)$ is denoted
by $\ell(\gamma)$.

\begin{definition}
	Let $\xi\in\partial\mathbb H^2$ and let $(\sigma(t))_{t\ge 0}$ be a geodesic ray
	converging to $\xi$. The \emph{Busemann cocycle based at $\xi$} is the function
	\[
	B_\xi(x,y):=\lim\limits_{t\to+\infty}\bigl(d(x,\sigma(t))-d(y,\sigma(t))\bigr),
	\qquad x,y\in\mathbb H^2.
	\]
\end{definition}

The Busemann cocycle satisfies the cocycle relation
\[
B_\xi(x,z)=B_\xi(x,y)+B_\xi(y,z)
\qquad \text{for all } x,y,z\in \mathbb H^2.
\]

\begin{definition}
	Fix $\xi\in\partial\mathbb H^2$ and $x_0\in\mathbb H^2$. A \emph{horocycle centered at $\xi$} passing through $x_0$
	is a level set of the map
	\[
	x\longmapsto B_\xi(x,x_0).
	\]
	Equivalently, in the upper half--plane model, horocycles are Euclidean circles tangent
	to $\partial\mathbb H^2=\mathbb R\cup\{\infty\}$ together with horizontal lines.
\end{definition}

We denote by $T^1\mathbb H^2$ the unit tangent bundle of $\mathbb H^2$.
For $\tilde v\in T^1\mathbb H^2$, we write $\tilde v(t)$ for the point at time $t$
along the geodesic determined by $\tilde v$.


%is the one-parameter family
%\[
%(g_t)_{t\in\mathbb R}:T^1\mathbb H^2\to T^1\mathbb H^2,
%\]
%defined by
%\[
%g_t(\tilde v)=\dot{\sigma}_{\tilde v}(t),
%\]
%where $\sigma_{\tilde v}$ is the geodesic with initial tangent vector $\tilde v$.

The \emph{horocyclic flow} on $T^1\mathbb H^2$ is the one-parameter family
\[
(h_s)_{s\in\mathbb R}:T^1\mathbb H^2\to T^1\mathbb H^2,
\]
where $h_s(\tilde v)$ is obtained by moving $\tilde v$ along the horocycle passing through its basepoint $\tilde v(0)$ and centered at 
its forward endpoint $\tilde v(+\infty)$, with parameter $s$ equal to signed
hyperbolic arc-length. The \emph{geodesic flow} on $T^1\mathbb H^2$ will be denoted by $
(g_t)_{t\in\mathbb R}:T^1\mathbb H^2\to T^1\mathbb H^2. 
$ 


Although there is a natural distance, called \emph{Sasaki distance} (see \cite{Sasaki}) on $T^1\H^2$, we will consider the following distance \cite{Ballmann}: for $\tilde v,\tilde w\in T^1\mathbb H^2$, define
\[
d_1(\tilde v,\tilde w):=d(\tilde v(0),\tilde w(0))+d(\tilde v(1),\tilde w(1)).
\]



From now on, we fix a torsion-free discrete subgroup
$
\Gamma<\operatorname{PSL}(2,\mathbb R),
$
and we consider the hyperbolic surface
$
S=\Gamma\backslash \mathbb H^2
$
and its unit tangent bundle
$
T^1S=\Gamma\backslash T^1\mathbb H^2.
$

We use a tilde to denote lifts to the universal cover: if $u\in T^1S$, then
$\tilde u\in T^1\mathbb H^2$ denotes any lift of $u$.

Since the geodesic flow and the horocyclic flow on $T^1\mathbb H^2$ commute with
the action of $\Gamma$, they descend to flows on $T^1S$, still denoted by
$
(g_t)_{t\in\mathbb R}
$ \text{and} $
(h_s)_{s\in\mathbb R}.
$
These are respectively called the \emph{geodesic flow} and the \emph{horocyclic flow}
on $T^1S$.


The distance $d_1$ in $T^1\H^2$ induces a distance on $T^1S$ that we will also denote by $d_1$.

\begin{definition}
	Let $u\in T^1S$. The \emph{strong-stable set} of $u$ is
	\[
	W^{ss}u:=\left\{v\in T^1S:\lim\limits_{t\to+\infty} d_1(g_t u,g_t v)=0\right\}.
	\]
\end{definition}

\begin{remark}
 In Appendix \ref{apendice}  we prove that the distances $d_1$ and $d_{S_A}$ define the same strong-stable set. 
\end{remark}

%As explained in Appendix~A, the distance $d_1$ is equivalent to the Sasaki metric,
%so the above definition is equivalent to the usual one.
	

		
Throughout this paper, $\tilde u_0$ in $T^1\H^2$ is defined by $\tilde u_0(0)=i$, $\tilde u_0(+\infty)=\infty$.
	
	\begin{definition}\label{nested_sequence}
%		A hyperbolic flute $S_F=\Gamma\backslash \H^2$ is hyperbolic surface where the Fuchsian group $\Gamma$ is generated by a sequence of hyperbolic isometries $(\gamma_n)_{n\in\N}$ satisfying:
%		\begin{enumerate}
%			\item $|\gamma_n^-|<|\gamma_{n+1}^-|$ and $|\gamma_n^+|<|\gamma_{n+1}^+|$.
%			\item $\lim\limits_{n\to \infty} \gamma_n^-=\lim\limits_{n\to \infty} \gamma_n^+=\infty.$ 
%			\item The signs of $\gamma_n^+$ and $\gamma_n^-$ are different for all $n$ and, for all $n$,  $\gamma_n$ has the same sign.
%			\item The half planes $\H^c_i(\gamma_n)$ are all disjoints. Furthermore, the sign of every extremety $e_1(\gamma_n)$ of the bisector $\partial \H_i (\gamma_n)$ is the same for all $n\in N$. The same applys for $e_2(\gamma_n)$. In addition, for every $n\in \N$: 
%			
%			$$0<|e_1(\gamma_n)|<|\gamma_n^+|<|e_2(\gamma_n)|<|e_1(\gamma_{n+1})|<|\gamma_{n+1}^+|.$$
%		\end{enumerate}
Let $(\gamma_n)_{n\geq 0}$ be a sequence of hyperbolic isometries. We say that  $(\gamma_n)_{n\geq 0}$ is a nested sequence if:
\begin{enumerate}
	\item The group $F=\langle \gamma_n \rangle$ is discrete,
	\item  $(\gamma_n^+)_{n\geq 0}$ is an increasing sequence of $\R^+$,
	\item  $(\gamma_n^-)_{n\geq 0}$ is a decreasing sequence of $\R^-$,
	\item $\forall n\geq 0$, $(\gamma_n^-,\gamma_n^+)\cap \Tilde u_0(\R^+)=\{ \tilde u_0(t_n) \}$ with $t_n>0$,
	\item $(\ell(\gamma_n))_{n\geq 0}$ is decreasing and convergent to $0$.
\end{enumerate}
	\end{definition}
	
	\begin{property}
		Let $(\gamma_n)_{n\geq 0}$ be a nested sequence. Then
		\begin{enumerate}
				\item $\lim\limits_{n\to+\infty}\gamma_n^+=\lim\limits_{n\to+\infty}\gamma_n^-=\infty$,
			\item $\lim\limits_{n\to \infty} t_n=+\infty$.
		
		\end{enumerate}
	\end{property}
	
	\begin{proof}
%		(1)Direct consequence of the Ping-Pong Dynamics (see \cite{Dalbo}).
%			(2),(3) $\Gamma$ acts properly discontinuosly on $\H^2$ and $\lim\limits_{n\to\infty}\ell(\gamma_n)=0$, then the increasing sequence $(t_n)_{n\in\N}$ cannot acumulate on $\H^2$ and thus $\lim\limits_{n\to\infty}t_n=+\infty$ and $\lim\limits_{n\to\infty}\gamma_n^+=\lim\limits_{n\to\infty}\gamma_n^-=+\infty$.

Let us see that
\[
\lim\limits_{n\to+\infty}\gamma_n^+=\lim\limits_{n\to+\infty}\gamma_n^-=\infty.
\]

Suppose that this is not the case. We may assume that
\[
\lim\limits_{n\to+\infty}\gamma_n^+=\xi\in\mathbb R^+\cup\{\infty\},
\qquad
\lim\limits_{n\to+\infty}\gamma_n^-=\eta\in\mathbb R^-\cup\{\infty\}.
\]


\medskip

\noindent
If $\xi\neq \infty$, or $\eta\neq \infty$,
then $\eta\neq \xi$, and the geodesics $(\gamma_n^-,\gamma_n^+)$ converge to the geodesic $(\eta,\xi)$.
Hence $i$ stays at bounded distance from $(\gamma_n^-,\gamma_n^+)$.
Since $\ell(\gamma_n)\to 0$, it follows that
$
d(\gamma_n i,i)
$ is bounded uniformly in $n$, which contradicts the discreteness of $F$.


Thus, $\eta=\xi=\infty$, and conclude that $\lim\limits_{n\to \infty} t_n=+\infty$ from the definition of $t_n$.
	\end{proof}	




\section{Winding around a closed geodesic}

\begin{definition}
	Let $\tilde{u}\in T^{1}\mathbb{H}^{2}$ and $\gamma$ a hyperbolic isometry such that
	\[
	\tilde{u}(\mathbb{R}^{+})\cap (\gamma^{-},\gamma^{+})\neq \varnothing.
	\]
	We define $\mathrm{Wind}_{\gamma}(\tilde{u})\in T^{1}\mathbb{H}^{2}$ by:
	\begin{itemize}
		\item $\mathrm{Wind}_{\gamma}(\tilde {u})(0)=\tilde{u}(0)$,
		\item $\mathrm{Wind}_{\gamma}(\tilde{u})(+\infty)=\gamma\bigl(\tilde{u}(+\infty)\bigr)$.
	\end{itemize}
\end{definition}

	Geometrically, when $\gamma$ belongs to a Fuchsian group $\Gamma$, on the surface  $S=\Gamma \backslash\H^2$, the projection of the ray $\mathrm{Wind}_{\gamma}(\tilde{u})(\mathbb{R}^{+})$, denoted by $\mathrm{Wind}_{\gamma}(u)(\mathbb{R}^{+})$, 
is obtained from ${u}(\mathbb{R}^{+})$ after winding around the closed geodesic
corresponding to $(\gamma^{-},\gamma^{+})$. %notation in the quotient...
 
	\begin{figure}[htbp]
	\centering
	\begin{normalsize} 
		%\def\svgwidth{\textwidth}
		\def\svgwidth{13cm}
		%\def\svgscale{0.7}
		\input{{figs/}dessin12.pdf_tex}
	\end{normalsize}
	\caption{Winding around a closed geodesic in the universal covering, where $t_0$ the time such that $\tilde u(t_0)\in (\gamma^-,\gamma^+)$.}
	\label{dessin12}
\end{figure}


\begin{definition}[winding time]\label{def_winding_time}

The winding time associated to $\operatorname{Wind}_\gamma(\tilde u)$ is the real number $\tau_{\gamma,\tilde u}$ defined by
$$
\tau_{\gamma,\tilde u}=B_{\tilde u(+\infty)}(\gamma^{-1}\tilde u(0), \tilde u(0)).
$$

\end{definition}

We clearly have $g_{\tau_{\gamma,\tilde u}}(\operatorname{Wind}_\gamma(\tilde u))\in h_{\R}( \gamma \tilde u)$.

 \begin{proposition}[bound of the winding time]\label{bound_wind}
 	Let $\tilde{u}\in T^{1}\mathbb{H}^{2}$ and $\gamma$ a hyperbolic isometry such that the geodesic ray $\tilde u(\R^+)$ meets the axis $(\gamma^-,\gamma^+)$, 
 	then
 	\[
 	\left| \tau_{\gamma,\tilde u} \right| \leq \ell(\gamma).
 	\]
 	
 \end{proposition}
 \begin{proof}
 	Let $p$ be the intersection point between $(\gamma^-,\gamma^+)$ and $\tilde u(\R^+)$.
 	As $p$ belongs to $\tilde u(\R^+)$, we have
 	$$
 	B_{\tilde u(+\infty)}(p,\tilde u(0))=-d(p,\tilde u(0)).
 	$$
 	In addition, as $p\in (\gamma^-,\gamma^+)$ we know that $d(\gamma^{-1}p,p)=\ell(\gamma)$. Thus
 	\begin{align*}
 		B_{\tilde u(+\infty)}(\gamma^{-1}\tilde u(0),\tilde u(0)) 
 		& =  B_{\tilde u(+\infty)}(\gamma^{-1}\tilde u(0),\gamma^{-1}p)+  
 		\\
 		&+B_{\tilde u(+\infty)}(\gamma^{-1}p,p) +B_{\tilde u(+\infty)}(p,\tilde u(0)) \\
 		& \leq  d(\gamma^{-1}\tilde u(0),\gamma^{-1}p)+\ell(\gamma)-d(p,\tilde u(0))\\
 		& = \ell(\gamma).
 	\end{align*}
 	
 	For the other inequality, we take $p_0$ in the intersection $[\gamma^{-1}\tilde u(0),\tilde u(+\infty) )\cap (\gamma^-,\gamma^+)$. This intersection is nonempty, since the isometry $\gamma^{-1}$ preserves the two connected components of $\H^2\setminus(\gamma^-,\gamma^+)$; hence $\gamma^{-1}\tilde u(0)$ lies in the same connected component as $\tilde u(0)$.
 	
 	Then $B_{\tilde u(+\infty)}(\gamma^{-1}\tilde u(0),p_0)=d(\gamma^{-1}\tilde u(0),p_0)$ and $d(p_0,\gamma p_0)=\ell(\gamma)$. Therefore 	
 	\begin{align*}
 		B_{\tilde u(+\infty)}(\gamma^{-1}\tilde u(0),\tilde u(0)) &= B_{\tilde u(+\infty)}(\gamma^{-1}\tilde u(0),p_0)+\\
 		&+B_{\tilde u(+\infty)}(p_0,\gamma p_0)+B_{\tilde u(+\infty) }(\gamma p_0,\tilde u(0)) \\
 		& \geq d(\gamma^{-1}\tilde u(0),p_0)-d(p_0,\gamma p_0)-d(\gamma p_0,\tilde u(0))\\
 		&= - \ell (\gamma).
 	\end{align*}
 \end{proof}
 
 The following proposition will be key for the proof of the main theorem.
 
 %$d_1\big(g_{t+\tau_{\gamma,\tilde u}} v, g_t u\big)$ in terms of the length of $\gamma$, where $v(\R^+)$ is a geodesic ray that winds once around $(\gamma^-,\gamma^+)$ with respect to $u(\R^+)$.
 
 \begin{proposition}[Key Proposition]\label{key_prop} 
 	Let $\tilde{u}\in T^{1}\mathbb{H}^{2}$ and $\gamma$ a hyperbolic isometry such that the geodesic ray $\tilde u(\R^+)$ meets the axis $(\gamma^-,\gamma^+)$. Then
 	\[\forall\, t\geq0, ~
 	\min \{ d_1\!\left(g_{t+\tau_{\gamma,\tilde u}}\bigl(\operatorname{Wind}_{\gamma}(\tilde u)\bigr),\, g_t\tilde u\right),
 	d_1\!\left(g_{t+\tau_{\gamma,\tilde u}}\bigl(\operatorname{Wind}_{\gamma}(\tilde u)\bigr),\, \gamma g_t\tilde u\right)
 	\}
 	\le 8 \ell(\gamma).
 	\]
 \end{proposition}
 \begin{proof}
 	%	We are going to proof that for all $t>0,$
 	%	\[
 	%	\Delta_t^{d}
 	%	:=\min\Bigl(
 	%	d\bigl(g_{t+\tau_{\gamma,\tilde u}}\,\mathrm{Wind}_{\gamma}(\tilde u),\,g_t\tilde u\bigr),
 	%	d\bigl(g_{t+\tau_{\gamma,\tilde u}}\,\mathrm{Wind}_{\gamma}(\tilde u),\,\gamma\,g_t\tilde u\bigr)
 	%	\Bigr)\le \varepsilon ,
 	%	\]
 	%	which implies the desired statement. 
 	For simplicity, we denote $\tilde v=\mathrm{Wind}_{\gamma}(\tilde u)$ and $\tau=\tau_{\gamma,\tilde u}$.
 	Up to applying an isometry of $\mathbb{H}^2$, we can assume that $\tilde u(0)=i$, $\tilde u(+\infty)=\infty$ and $\gamma^{+}>0$.
 	
 	Let $t_\gamma\ge 0$ be such that $\tilde u(t_\gamma)\in(\gamma^{-},\gamma^{+})$. Since $i$ and $\gamma \, \infty$ lie in different half-planes bounded by $(\gamma^{-},\gamma^{+})$, the geodesic ray $\tilde v(\mathbb R^{+})$ must intersect the axis $(\gamma^{-},\gamma^{+})$. Let $t'_\gamma\ge 0$ be such that $\tilde v(t'_\gamma)\in(\gamma^{-},\gamma^{+})$.
 	
 	We show that the intersection point $\tilde v(t'_\gamma)\in(\gamma^{-},\gamma^{+})$ lies in the following hyperbolic segment:
 	\[\tilde v(t'_\gamma)\in
 	[\tilde{u}(t_\gamma), \gamma \tilde{u}(t_\gamma)].
 	\]
 	We observe that $\tilde v(+\infty)$ is the real interval $]\gamma^+, +\infty[$ and $\tilde v(-\infty)$ is in the interval $]\gamma^-,0[$. Hence, $\tilde v(t'_\gamma)$ is on the right of $\tilde u(t_\gamma)$.  Applying $\gamma^{-1}$, we obtain that $\gamma^{-1} \tilde v (+\infty) = \infty$ and $\gamma^{-1} \tilde v (-\infty)$ is also on the interval $]\gamma^-,0[$. As a consequence, $\gamma^{-1} \tilde v(t'_\gamma)$, which is the intersection of the vertical ray $\gamma^{-1}\tilde{v}(\mathbb{R}^+)$ and the geodesic $(\gamma^{-}, \gamma^{+})$, is at left of the point $\tilde{u}(t_\gamma)$.
 	Therefore, $\tilde v(t'_\gamma)\in(\gamma^{-},\gamma^{+})$ has to be on the geodesic segment
 	$[\tilde{u}(t_\gamma), \gamma \tilde{u}(t_\gamma)]$. %picture
 	
 	We conclude that $d(\tilde v(t'_\gamma),\tilde{u}(t_\gamma) )$ and $d(v(t'_\gamma), \gamma \tilde{u}(t_\gamma) )$ are both less than $\ell(\gamma)$. Moreover, by a triangular inequality,
 	\[
 	|t_\gamma - t'_\gamma | \le d(\tilde v(t'_\gamma),\tilde{u}(t_\gamma) ) \le \ell(\gamma).
 	\]
 	
% 	We shall remark that, because
% 	\[
% 	d(i, (\gamma^{-}, \gamma^{+})) \ge 1,
% 	\]
% 	we have
% 	\[
% 	d(i, \tilde{u}(t_\gamma)) \ge 1.
% 	\]
% 	In addition, we showed that
% 	\[
% 	\tilde{v}(\mathbb{R}^+) \cap
% 	[\tilde{u}(t_\gamma), \gamma \tilde{u}(t_\gamma)]
% 	\supset \{\tilde{v}(t'_\gamma)\},
% 	\]
% 	so
% 	\[
% 	d(\tilde{u}(t_\gamma), \tilde{v}(t'_\gamma)) \le \ell(\gamma),
% 	\]
% 	which means we can choose the length $\ell(\gamma)$ to be arbitrarily
% 	small so that $\tilde{v}(\mathbb{R}^+)$ is arbitrarily near
% 	$\tilde{u}(\mathbb{R}^+)$.
% 	In other words, the smaller $\ell(\gamma)$, the bigger $\tilde{v}(+\infty)$.
% 	
% 	\begin{lemma}
% 		\begin{enumerate}
% 			\item[(i)] $t_\gamma\geq  1$,
% 			\item[(ii)] $d\bigl(\tilde u(t_\gamma),\tilde v(t'_\gamma)\bigr)\le \ell(\gamma)$,
% 			\item[(iii)] $|t_\gamma-t'_\gamma|\le \ell(\gamma)$.
% 		\end{enumerate}
% 	\end{lemma}
% 	
% 	\begin{proof}[proof of the lemma]
% 		\begin{enumerate}
% 			\item[(i)] $d\bigl(i,\tilde u(t_\gamma)\bigr)\ge d\bigl(i,(\gamma^{-},\gamma^{+})\bigr)\ge 1$.
% 			\item[(ii)] The geodesic ray $\gamma^{-1}\tilde v(\mathbb R^{+})$ is vertical; its origin is also the origin of $\gamma^{-1}\tilde u(\mathbb R^{+})$. 
% 			$$
% 			\gamma^{-1}\tilde v(t'_\gamma)\in [\gamma^{-1}\tilde u (t_\gamma), \tilde u(t_\gamma)],
% 			$$
% 			as a consequence $d(\tilde u(t_\gamma),\tilde v(t_\gamma'))\leq \ell(\gamma)$.
% 			\item[(iii)] By triangle inequality,
% 			\[
% 			\bigl|\,d\bigl(i,\tilde v(t'_\gamma)\bigr)-d\bigl(i,\gamma^{-1}\tilde v(t'_\gamma)\bigr)\,\bigr|
% 			\le d\bigl(\tilde u(t_\gamma),\tilde v(t'_\gamma)\bigr),
% 			\]
% 			and applying (ii) yields (iii).
% 		\end{enumerate}
% 	\end{proof}
% 	
 	
 	% -------------------- Continuation (keyprop2) --------------------
 	
 	We first prove Proposition \ref{key_prop} on $S$, replacing $d_1$ by $d$.	
 	
% 	We introduce for all $t\ge 0,$
% 	\[
% 	\Delta_t^{d}
% 	:=\min\Bigl(d\bigl(\tilde v(t+\tau),\tilde u(t)\bigr),\;
% 	d\bigl(\tilde v(t+\tau),\gamma\,\tilde u(t)\bigr)\Bigr).
% 	\]
% 	We prove
% 	\[
% 	\forall\,t\ge 0,\qquad \Delta_t^{d}\le 3\ell(\gamma).
% 	\]
 	
 	\noindent\textbf{Case 1.}  If $0\le t\le t_\gamma$, then 
 	\[
 	d\bigl(\tilde v(t),\tilde u(t)\bigr)\le 2\ell(\gamma).
 	\] 
% 	By the triangular inequality and Proposition \ref{bound_wind}, we have
% 	\begin{align*}
% 		d\bigl(\tilde v(t+\tau),\tilde u(t)\bigr)
% 		&\le  d\bigl(\tilde v(t+\tau),\tilde v(t)\bigr)+d\bigl(\tilde v(t),\tilde u(t)\bigr)\\
% 		&\le |\tau|+d\bigl(\tilde v(t),\tilde u(t)\bigr)\\
% 		&\le \ell(\gamma)+d\bigl(\tilde v(t),\tilde u(t)\bigr).
% 	\end{align*}

 Since the distance in $\mathbb{H}^2$ between two geodesic rays starting at the same point is increasing, it follows 
 \[
  	\forall\,0\le t\le t_\gamma,\qquad
  	d\bigl(\tilde v(t),\tilde u(t)\bigr)\le d\bigl(\tilde v(t_\gamma),\tilde u(t_\gamma)\bigr).
  	\]
 %I think that this is a well known fact and we can avoid the proof: 
% 	The triangle $[\,i,\tilde u(t),\tilde v(t)\,]$ is ; denote by $\beta$ the angle at $i$ and by $\alpha_t$ the angle at $\tilde u(t)$ (or $\tilde v(t)$).
% 	By the hyperbolic sine formula (\cite{Beardon} \S7.12),
% 	\[
% 	\sinh d\bigl(\tilde u(t),\tilde v(t)\bigr)=\frac{\sinh t\;\sin\beta}{\sin\alpha_t}.
% 	\]
% 	Since $\pi-\beta-2\alpha_t$  is increasing and $\alpha_t\to 0$ as $t\to 0$, it follows that
% 	\[
% 	\forall\,0\le t\le t_\gamma,\qquad
% 	d\bigl(\tilde u(t),\tilde v(t)\bigr)\le d\bigl(\tilde u(t_\gamma),\tilde v(t_\gamma)\bigr).
% 	\]
Using the bounds established above, one obtains
 	\[
 	d\bigl(\tilde v(t_\gamma),\tilde u(t_\gamma)\bigr) \le 
 	|t_\gamma - t'_\gamma | + d\bigl(\tilde v(t'_\gamma),\tilde u(t_\gamma)\bigr)  \le
 	2\ell(\gamma),
  	\]
which implies the statement.
 	
% 	We have shown 
%\[
%\forall\,0\le t\le t_\gamma,\qquad
%d\bigl(\tilde v(t+\tau),\tilde u(t)\bigr)\le 3\ell(\gamma).
%\] 
% 
 	
 	\smallskip
 	\noindent\textbf{Case 2:} If $ t\ge  t_\gamma$, then 
 	\[
 	d\bigl(\tilde v(t), \gamma \tilde u(t)\bigr)\le 2\ell(\gamma).
 	\] 
 	
% 	Similarly as above, 
% 	\begin{align*}
% 		d\bigl(\tilde v(t+\tau), \gamma \tilde u(t)\bigr) &= d\bigl(\gamma^{-1}\tilde v(t+\tau),\tilde u(t)\bigr) \\
% 		&\le d\bigl(\gamma^{-1}\tilde v(t+\tau),\gamma^{-1}\tilde v(t)\bigr)
% 		+d\bigl(\gamma^{-1}\tilde v(t),\tilde u(t)\bigr)\\
% 		&\le \ell(\gamma)+d\bigl(\gamma^{-1}\tilde v(t),\tilde u(t)\bigr).
% 	\end{align*}
 	Put
 	\[
 	\gamma^{-1}\tilde v(t)=a+i\,e^{\tau+t}\quad\text{and}\quad \tilde u(t)=i\,e^{t}.
 	\]
 	We have \cite[\S7.20]{Beardon} 
 	\begin{equation*}%\tag{$\ast$}
 		\left(\sinh\frac{d\bigl(\gamma^{-1}\tilde v(t),\tilde u(t)\bigr)}{2}\right)^{2}
 		=
 		\frac{a^{2}+e^{2t}(e^{\tau}-1)^{2}}{4e^{\tau}e^{2t}}
 		=
 		\frac{a^{2}}{4e^{\tau}}\,e^{-2t}
 		+\frac{(e^{\tau}-1)^{2}}{4e^{\tau}}.
 	\end{equation*}
 	It follows that $d\bigl(\gamma^{-1}\tilde v(t),\tilde u(t)\bigr)$ is decreasing, and hence for $t\ge t_\gamma$,
 	\[
 	d\bigl(\gamma^{-1}\tilde v(t),\tilde u(t)\bigr)
 	\le d\bigl(\gamma^{-1}\tilde v(t_\gamma),\tilde u(t_\gamma)\bigr).
 	\]
 	Moreover,
 	\[
 	d\bigl(\gamma^{-1}\tilde v(t_\gamma),\tilde u(t_\gamma)\bigr)
 	\le |t_\gamma-t'_\gamma| + d\bigl(\gamma^{-1}\tilde v(t'_\gamma), \tilde u(t_\gamma) \bigr).
 	\]
 	By the bounds proved above, each term on the right hand side is bounded by $\ell (\gamma)$.
 	We conclude that, for $t\ge t_\gamma$, $d\bigl(\tilde v(t), \gamma \tilde u(t)\bigr) = d\bigl( \gamma^{-1} \tilde v(t), \tilde u(t)\bigr)  \le 2\ell(\gamma)$. 
 	
 	\smallskip
 	
 Let us now prove Proposition \ref{key_prop}. Recall that 
 \[
d_1\!\left(g_{t} \tilde v ,\, g_t\tilde u\right) = d(\tilde v (t), \tilde u (t)) + d(\tilde v (t+1), \tilde u (t+1)). 
 \]
 We distinguish three scenarios:
 
 \begin{enumerate}[(a)]
 	\item If $t \le t_\gamma -1 $, by the first case, we have that both
 	$d(\tilde v(t), \tilde u(t))$ and \\
 	$d(\tilde v(t+1), \tilde u(t+1))$
 	are bounded by $2\ell(\gamma)$. Then
 	$
 	d_1(g_{t}\tilde v, g_t \tilde u) \le 4\ell(\gamma).
 	$
 	
 	\item If $t \ge t_\gamma$, by the second case, we can still
 	bound both
 	$d(\tilde v(t), \gamma \tilde u(t))
 	$ and $ 
 	d(\tilde v(t+1), \gamma \tilde u(t+1))$
 	by $2\ell(\gamma)$. Then
 	$
 	d_1(g_{t}\tilde v, \gamma g_t \tilde u) \le 4 \ell(\gamma).
 	$
 	
 	\item If $t_\gamma -1 < t < t_\gamma$. We have
 	\begin{align*}
 		d_1(g_{t}\tilde v, g_t \tilde u) &= d(\tilde v (t), \tilde u (t)) + d(\tilde v (t+1), \tilde u (t+1)) \\
 		&\le  d(\tilde v (t), \tilde u (t)) + d(\tilde v (t+1), \gamma \tilde u (t+1)) + d( \gamma \tilde u (t+1), \tilde u (t+1)).
 	\end{align*}
	The first and the second terms in the last line are bounded by $2\ell (\gamma)$ by Cases 1 and 2 above, respectively. 
 	
	In order to bound the third term we apply the formula of the displacement function \cite[Theorem 7.35.1]{Beardon}. Writing $z= \tilde u(t+1)$, we obtain
 	\begin{equation*}\label{eq:5.11}
 		\sinh\!\left(\frac{d(z, \gamma z)}{2}\right)
 		=
 		\cosh \bigl(d(z, (\gamma^-, \gamma^+))\bigr)
 		\sinh\!\left(\frac{\ell(\gamma)}{2}\right).
 	\end{equation*}
 
 We notice that \[ d(z, (\gamma^-, \gamma^+) )  \le d(\tilde u(t+1), \tilde u (t_\gamma) ) = t+1 - t_\gamma \le 1 .\] We obtain 
 \[
 \sinh\!\left(\frac{d(z, \gamma z)}{2}\right) \leq 2 \sinh\!\left(\frac{\ell(\gamma)}{2}\right).
 \]
 Since $2\sinh(\frac{\ell(\gamma)}{2})\leq \sinh \ell(\gamma)$, then the third term $d(z,\gamma z)$ is bounded by $2\ell(\gamma)$. We obtain that, if $t_\gamma -1 < t < t_\gamma$, then 
 \[
 d_1(g_{t}\tilde v, g_t \tilde u) \le 6 \ell(\gamma).
 \]
 
 
 \end{enumerate}

Finally, we observe that 
\begin{align*}
	d_1\!\left(g_{t+\tau} \tilde v ,\, g_t\tilde u\right) & \le d_1\!\left(g_{t+\tau} \tilde v ,\, g_t\tilde v\right) + d_1\!\left(g_{t} \tilde v ,\, g_t\tilde u\right)\\
	& =  2|\tau| + d_1\!\left(g_{t} \tilde v ,\, g_t\tilde u\right) \\
	&\le  2\ell(\gamma) + d_1\!\left(g_{t} \tilde v ,\, g_t\tilde u\right),
\end{align*}
and similarly we have
\[
d_1\!\left(g_{t+\tau} \tilde v ,\, \gamma g_t\tilde u\right) \le 2\ell(\gamma) + d_1\!\left(g_{t} \tilde v ,\, \gamma g_t\tilde u\right).
\]
Together with the previous bounds, these imply the statement.
 \end{proof}

%In order to prove the aforementioned theorems, several preparatory lemmas will be required. Additionally, for simplification purposes, we will assume that  $\tilde u(0)=i$ and that $\tilde u(+\infty)=\infty$, which is always true up to conjugation of the group $\Gamma$. The first lemma will facilitate the selection of a specific sequence of elements within our group $\Gamma$. We recall that the notation $\H_i^c(\gamma_n)$ denotes the complementary of the half plane $\H_i(\gamma_n)=\{z\in\H^2:~ d(z,i)\leq d(z,\gamma(i))\}$.  


\section{Proof of Main Theorem}
\begin{proposition}\label{prop3}
	Let $S=\Gamma\backslash\H^2$ be a hyperbolic surface and $u\in T^1S$. If $u(\R^+)$ crosses infinitely many closed geodesics with length converging to $0$, then there exists a sequence $(\gamma_n')_{n\geq 0}$ of hyperbolic isometries in $\Gamma$, and an isometry $f$ such that:
	\begin{enumerate}
		\item $(\gamma_n=f\gamma_n'f^{-1})_{n\geq 0}$ is a nested sequence (Definition \ref{nested_sequence}),
		\item The geodesic ray $f^{-1}(\tilde u_0)(\R^+)\subset \H^2$ projects onto $u(\R^+)\subset S=\Gamma\backslash\H^2$.
	\end{enumerate}
\end{proposition}

\begin{proof}
	Let $\tilde u\in T^1\mathbb H^2$ be a lift of $u$.	
	Since $\mathrm{PSL}(2,\mathbb R)$ acts transitively on $T^1\mathbb H^2$, there exists an isometry
	$
	f\in \mathrm{PSL}(2,\mathbb R)
	$
	such that
	$
	\tilde u=f^{-1}(\tilde u_0).
	$
	Therefore the geodesic ray
	$
	f^{-1}(\tilde u_0)(\mathbb R^+)=\tilde u(\mathbb R^+)
	$
	projects onto
	$
	u(\mathbb R^+)\subset S.
	$
	
	By assumption, the ray $u(\mathbb R^+)$ crosses infinitely many closed geodesics whose lengths converge to $0$. For each such closed geodesic, choose a hyperbolic element
	$
	\gamma_n'\in \Gamma
	$
	whose axis projects onto that closed geodesic and meets the ray $\tilde u(\R^+)$. Passing to a subsequence, we may assume that
	$
	\ell(\gamma_n')
	$
	is decreasing and converges to $0$.
	
	Set
	$
	\gamma_n=f\gamma_n'f^{-1}.
	$
	Clearly, each $\gamma_n$ is hyperbolic and
	$
	\ell(\gamma_n)=\ell(\gamma_n').
	$
	Moreover, since the axis of $\gamma_n'$ meets $\tilde u(\mathbb R^+)$, the axis of $\gamma_n$ meets
	$
	f(\tilde u(\mathbb R^+))=\tilde u_0(\mathbb R^+).
	$
	
	We shall now prove that, after possibly replacing some $\gamma_n$ by their inverses and extracting a subsequence, $(\gamma_n)_{n\geq 0}$ is a nested sequence.
	
	\medskip
	
	\noindent
	Let
	$
	F=\langle \gamma_n \; ;\; n\geq 0\rangle.
	$
	Since
	$
	F\subset f\Gamma f^{-1}
	$
	and $\Gamma$ is discrete, $F$ is discrete. Thus, condition~(1) is verified.
	
	\medskip
	
	\noindent
	Since the axis of each $\gamma_n$ intersects the vertical ray
	$
	\tilde u_0(\mathbb R^+)=[i,\infty),
	$
	its two endpoints lie on opposite sides of $0$. Hence, after possibly replacing $\gamma_n$ by $\gamma_n^{-1}$, we may assume that
	$
	\gamma_n^-<0<\gamma_n^+
	 \text{ for all }n,
	$ and that the sequence $(\gamma_n)_{n\geq 0}$ and $(\gamma_n^-)_{n\geq 0}$ are increasing and decreasing, respectively, which gives conditions~(2) and (3). 
	
	By construction, the axis of each $\gamma_n$ intersects $\tilde u_0(\mathbb R^+)$. Since this axis has endpoints $\gamma_n^-<0<\gamma_n^+$, the intersection consists of exactly one point. Therefore, for every $n\geq 0$, there exists $t_n>0$ such that
	\[
	(\gamma_n^-,\gamma_n^+)\cap \tilde u_0(\mathbb R^+)=\tilde u_0(t_n).
	\]
	This is condition~(4).
	
	\medskip
	
	\noindent
	Finally, since the sequence $(\ell(\gamma_n'))$ is chosen decreasing with limit $0$, and
	$
	\ell(\gamma_n)=\ell(\gamma_n'),
	$
	therefore $(\ell(\gamma_n))_{n\geq 0}$ is decreasing and converges to $0$. This is condition~(5).
	
	\medskip
\end{proof}




%	\begin{remark}
	%		The group $\Gamma_0$ of the proof of the Main Theorem in the case of a flute Surface will be the group generated by the sequence $(\gamma_n)_{n\in\N}$.
	%	\end{remark}
%	We now conclude by appyling the results of the previous section to the sequence $(\gamma_n)$.

Let $S=\Gamma\backslash \H^2$ and $u\in T^1 S$ satisfying the conditions of the Main Theorem \ref{thE} and $f\in\operatorname{PSL(2,\R)}$ given by Proposition \ref{prop3}, Clearly, it is enough to prove the Main Theorem for $S_0=f\Gamma f^{-1}\backslash \H^2$ and $u_0\in T^1S_0$ lifting to $\tilde u_0$.

Let $(\gamma_n')_{n\geq 0}$ given by Proposition \ref{prop3}, and set $\gamma_n=f\gamma_n'f^{-1}$, which is a nested sequence. %and $\Gamma_0=\langle \gamma_n\rangle$.
 Our goal is to give conditions on a subsequence $\alpha=(\alpha_n)_{n\geq0}$ of $(\gamma_n)_{n\geq 0}$ to guarantee the existence of $\tilde w_\alpha \in T^1\H^2$ with $\tilde w_\alpha(+\infty)=\lim\limits_{n\to +\infty}\alpha_0\ldots\alpha_n(\infty)$, such that its projection $w_\alpha$ onto $T^1S_0$ satisfies $ w_\alpha\in W^{ss}u_0-h_\R u_0$.


 
%  Observe that these assumptions entail no loss of generality, since one can always conjugate the underlying group $S_F$ so that the limit point of $\tilde u_0$ is $\infty$. Moreover, for any other $\tilde u(0)$ there exist $s_0,t_0\in \R$ such that $\tilde u(0)=g_{t_0}h_{s_0}(\tilde u(0))$. Therefore, the results of this section remain valid for the strong stable leaf of a vector $u$ that does not satisfy these two conditions.

%Our purpose is to construct vector $w_\alpha$ pointing towards $\lim\limits_{n\to\infty}\alpha_1\ldots\alpha_n(i)$
%satisfying the desired property, which is $w_\alpha\in W^{ss}u-h_\R u$ (this vector winds about each of the geodesics 
%$(\alpha_n^-,\alpha_n^+)$). 

%The position of the half planes $\H_i^c(\gamma_n)$ are illustrated in Figure \ref{dessin3}




\subsection{Definition of $v_n$ and $w_\alpha$}

\par\,

%In order to define our vector $w$ of the theorem \ref{toytheorem} we need the following Lemma, which will help us associate a sequence of vectors $(\tilde v_n)_{n\in \N}$ towards a vector $\tilde v\in T^1 \H^2$. 
	Let $(\alpha_n)_{n\in\\N}$ be a subsequence of $(\gamma_n)_{n\in \N}.$ We introduce $\beta_n=\alpha_0\ldots\alpha_n\in\Gamma_0,$ and $\tilde v_n\in T^1\H^2$ defined by:
	\begin{itemize}
		\item $\tilde v_n(0)=i,$
		\item $\tilde v_n(+\infty)=\beta_n(\infty)$.
	\end{itemize}
	
%	Note $\tilde v_n$ the vectors of $T^1 \H^2$ based on $i$ and directed towards $\beta_n(\infty), $ where $\beta_n=\alpha_1\ldots\alpha_n.$

\begin{lemma}\label{lemma_vectors}
	The sequence $(\tilde v_n)_{n\in\N}$ converges towards $\tilde v_\alpha\in T^1 \H^2 $ defined by $\tilde v_\alpha (0)=i$ and $\tilde v_\alpha(+\infty)=\lim\limits_{n\to +\infty}\beta_n(\infty)$.
	
\end{lemma}


\begin{proof}
	We have to prove that the sequence $(\beta_n(\infty))_{n\in\N}\subset \partial \H^2$ converges. 
	Using the dynamics of $\gamma_n$ we have:
	\begin{itemize}
		\item $\forall n\geq 1$, $0<\alpha_{n-1}\alpha_n\infty<\alpha_{n-1}\infty$,
		\item $\forall i\geq 0,$ if $0<x<y$, then $0<\alpha_i x<\alpha_iy.$
	\end{itemize}
	It follows that $(\beta_n(\infty))_{n\in\N}$ is a decreasing sequence of positive numbers and hence converges.
	
%	Let us prove that $\tilde v_n(\infty)=\beta_n(\infty)$ converge in $\H^{2}(\infty)$. As every $\tilde v_n$ is based in $i$, this proves convergence in $T^1\H^2$. To do so, we will show that the sequence $|\beta_n(\infty)|$ is decreasing and bounded from below by $|\alpha_0^+|$. 
%	Let $n$ be a non-negative integer and observe that, from the definition of flute (see \ref{def_flute}) and the dynamics of the action of the isometries of $\Gamma$ we have the following inequalities, as illustrated in Figure \ref{dessin14}.
	
%	\begin{figure}[htbp]
%		\centering
%		\begin{normalsize} 
%			%\def\svgwidth{\textwidth}
%			\def\svgwidth{13cm}
%			%\def\svgscale{0.7}
%			\input{{figs/}dessin14.pdf_tex}
%		\end{normalsize}
%		\caption{Position of the points $\alpha\cdot \infty$, $\alpha_{n+1}^+$ and $\alpha_{n+1}^-$ in $S^1$.}
%		\label{dessin14}
%	\end{figure}
	
%	\begin{equation}\label{eq13}
%	|\alpha_{n+1}^-|<|\alpha_{n+1}^+|<|\alpha_{n+1}\cdot\infty|<|\infty|,
%	\end{equation}
%	and also: 
%	$$0<|\alpha_0^+|<|\alpha_1^+|<\ldots<|\alpha_n^+|<|\alpha_{n+1}^+|.$$
%	
%	Now, if we keep applying the isometry $\alpha_k$ from $k=0$ to $k=n-1$ to the points $\alpha_{k+1}\ldots\alpha_{n+1}\cdot\infty$ and $\alpha_{k+1}\ldots\alpha_{n}\cdot\infty$ the same that we did in \ref{eq13} we obtain the following inequality (see Figure \ref{dessin15})
	
%	\begin{figure}[htbp]
%		\centering
%		\begin{normalsize} 
%			%\def\svgwidth{\textwidth}
%			\def\svgwidth{13cm}
%			%\def\svgscale{0.7}
%			\input{{figs/}dessin15.pdf_tex}
%		\end{normalsize}
%		\caption{Disposition of the limit points $\beta_n\cdot \infty$ and $\beta_{n+1}\cdot \infty$ in $S^1$.}
%		\label{dessin15}
%	\end{figure}
	
%	$$
%	|\alpha_0^+|<|\alpha_0\ldots\alpha_n\cdot\infty|<|\alpha_0\ldots\alpha_{n-1}\cdot\infty|<\ldots <|\alpha_0\cdot \infty|.
%	$$
%	Thus, for all $n$ we have that
%	
%	$$|\alpha_0^+|<|\beta_n\cdot\infty|<|\beta_{n-1}\cdot\infty|,$$
%	which concludes the proof. 

 \end{proof}


%PROPERTY2
%\begin{property}\label{property2}
%	$\tilde v_\alpha({+\infty})\notin\Gamma\infty$.
%\end{property}
%\begin{proof}
%	Suppose $\tilde v_\alpha(+\infty)=\gamma\infty $ with $\gamma\in \Gamma$ and $\gamma\neq Id$. Write $\gamma$ as a reduce word: $\gamma=\omega_0\ldots \omega_k$, with $\omega_i\in \{\gamma_n^{\pm 1},n\in \N\},$ $\omega_{i+1}\neq \omega_i^{-1}$, for $i\in \{0,\ldots,k-1\}$.
%	By construction $\tilde v_\alpha(+\infty)\in D(\alpha_0)$. Moreover, $\gamma\infty \in D(\omega_0)$, hence $\omega_0=\alpha_0$. By induction we obtain $\gamma=\alpha_0\ldots\alpha_k$. 
%	It follows that $$\infty=\gamma^{-1}\tilde v_\alpha(+\infty)\in D(\alpha_{k+1}),$$ which is a contradiction.
%\end{proof}

%	Suppose that there exists an isometry $\gamma\in\Gamma$ such that $\beta_\infty=\tilde v_\alpha(+\infty)= \lim\limits_{n\to\infty}\gamma_1\ldots\gamma_n(\infty)=\infty$. Observe that $\gamma=\gamma_{i_1}\ldots\gamma_{i_n}$ for certain $\gamma_{i_j}\in \Gamma$. Thus:
%	$$
%	\gamma_1\ldots\gamma_n\gamma_{n+1}\ldots...(\infty)=\gamma_{i_1}\ldots\gamma_{i_n}.
%	$$
%	But this implies that $\gamma_{i_1}\ldots\gamma_{i_n}=\gamma_1 \ldots\gamma_n$ and then
%	$$
%	\gamma_{n+1}\ldots (\infty)=\infty,
%	$$
%	which is contradiction because $\infty$ does not belong to $D(\gamma_{n+1})$.

Applying $\alpha_0^{-1}$ to $\tilde v_n(0)=i$ maps it into the same connected component of $\H^2-(\alpha_0^-,\alpha_0^+)$, repeating this argument inductively leads to the conclusion that $\alpha_{n}^{-1}\ldots\alpha_0^{-1}(i)$ belongs to the connected component of $\H^2-(\alpha_n^-,\alpha_n^+)$ containing $i$. It follows that
 the ray $\alpha_{n}^{-1}\ldots\alpha_0^{-1}\tilde v_n(\R ^+)=\beta_n^{-1}\tilde v_n(\R ^+)$ crosses $(\alpha_{n+1}^-,\alpha_{n+1}^+) $ and hence that $\tilde v_n(\R^+)$ crosses $(\beta_n(\alpha_{n+1}^-)),\beta_n(\alpha_{n+1}^+) )$.

As a consequence we have for all $n\geq 0$
$$
\tilde v_{n+1}=\mathrm{Wind}_{\beta_n\alpha_{n+1}\beta_n^{-1}}(\tilde v_n),
$$
and accordingly to Proposition \ref{bound_wind}
$$
\left|\tau_{\beta_n\alpha_{n+1}\beta_n^{-1},\tilde v_n}\right |\leq \ell(\alpha_{n+1}).
$$

	\begin{figure}[htbp]
	\centering
	\begin{normalsize} 
		%\def\svgwidth{\textwidth}
		\def\svgwidth{13cm}
		%\def\svgscale{0.7}
		\input{{figs/}dessin13.pdf_tex}
	\end{normalsize}
	\caption{$g_{r_n}\beta_n^{-1}\tilde v_n\in h_\R\tilde u$.}
	\label{dessin13}
\end{figure}


The next step of the construction is to ensure that the time spent to wind around all the geodesics $(\alpha_n^-,\alpha_n^+)$ is finite.

\begin{proposition}[Convergence of the time to wind around closed geodesics]\label{lema_4}
	
Let $r_n=B_\infty(\beta_n^{-1}i,i))$. If for every $n\geq 0$,  $\ell(\alpha_n)<\frac{1}{2^n}$, then the sequence $(r_n)_{n\in\N}$ converges towards a real number $r_\alpha$.    
\end{proposition}

\begin{proof}
We have $\left | r_{n+1}-r_n  \right |=\left | B_\infty(\alpha_{n+1}^{-1}\beta_n^{-1}(i), \beta_n^{-1}(i)) \right |$.
Hence
 $$\left | r_{n+1}-r_n  \right |=\left | B_{\beta_n \infty}(\beta_n\alpha_{n+1}^{-1}\beta_n^{-1}(i),i) \right |.$$
 It follows that

$$
\left | r_{n+1}-r_n  \right |=\left | \tau_{\beta_n\alpha_{n+1}\beta_n^{-1}}(\tilde v_n) \right|\leq \ell(\alpha_{n+1})\leq \frac{1}{2^{n+1}}.$$
As a consequence $(r_n)_{n\in \N}$ is a Cauchy sequence and hence converges.
\end{proof}

For a subsequence $\alpha$ of $(\gamma_n)_{n\N}$ satisfying the assumptions of Proposition \ref{lema_4} we will define $w_\alpha$ as follows:
$$w_\alpha=g_{r_\alpha}v_\alpha=\lim\limits_{n\to \infty}g_{r_n}v_n.$$

\begin{remark}\label{obs_orbit}
As direct consequence of the definition of $w_\alpha$ we deduce that $w_\alpha\in \overline{h_\R u_0}$.
\end{remark}






%
%\begin{proposition}[Key Proposition]\label{prop_retardo}
%	Given $\eta>0$, there exists $\varepsilon>0$ such that, for all hyperbolic isometry $\gamma$  verifying  $\ell(\gamma)<\varepsilon$ and, for $\tilde u, \tilde v\in T^1 \H^2$, the triple $(\tilde u, \tilde v, \gamma)$ is disposed in such a way that
%	\begin{enumerate}
%		\item $\tilde u(\R^+)$ meets the axis $(\gamma^-,\gamma^+)$,
%		\item  $\tilde u(0)=\tilde v(0)$,
%		\item $\tilde v(+\infty)=\gamma \tilde u(+\infty)$,
%		\item $d(\tilde u(0), (\gamma^{-}, \gamma^{+}))\geq 1$,
%	\end{enumerate}
%	then for all $t>0$:
%	
%	$$\Delta_t:=\min\{d_1(g_{t+r}\tilde v, g_t\tilde u), d_1(g_{t+r}\tilde v, \gamma g_t\tilde u)\}<\eta,
%	$$
%	
%	where $r=\text{B}_{\tilde u(+\infty)}(\gamma^{-1}\tilde u (0),\tilde u(0))$.
%\end{proposition}
%\begin{proof}
%	Up to conjugation in $\operatorname{PSL}_2(\mathbb{R})$ we can assume $\tilde u(0)=i,\, \tilde u(+\infty)=\infty$.
%	We will suppose $\gamma^+$ positive. If $\gamma^+>0,$ then $\gamma^-$ is negative and $\tilde v(+\infty)=\gamma \infty\geq\gamma^+$. If $\gamma^+$ is negative, the proof is symetric. 
%	
%
%		\begin{figure}[htbp]
%		\centering
%		\begin{normalsize} 
%			\def\svgwidth{\textwidth}
%			%\def\svgwidth{20cm}
%			%\def\svgscale{0.5}
%			\input{{figs/}dessin18.pdf_tex}
%		\end{normalsize}
%		\caption{The ray $\gamma^{-1}\tilde v(\R^+)$ meets the segment $[\gamma^{-1}\tilde u(t_\gamma),\tilde u(t_\gamma)]$.}
%		\label{dessin18}
%	\end{figure}
%	
%	
%	Let us begin by showing that the geodesic ray $\tilde v (\mathbb{R}^+)$ meets the hyperbolic segment $[\tilde u(t_\gamma), \gamma \tilde u(t_\gamma)]$, where $t_\gamma$ is the time such that $\tilde u(t_\gamma)\in (\gamma^-,\gamma^+)$.
%
%	Observe that the situation that we have is indeed the one in the Figure \ref{dessin6} because $\gamma^{-1}i$ lies at left of the vertical line $[0,\infty]$ as so does $\gamma^{-1}0$ because $\gamma^-<0<\gamma^+$. As a consequence of this, $\gamma^{-1}\tilde v(\R^+)$ is a vertical line that meets the geodesic $(\gamma^-,\gamma^+)$ at left of the point $\tilde u(t_\gamma)$ (see Figure \ref{dessin18}).
%	 Therefore, $\tilde v(\mathbb{R}^+)$ has to meet the segment $[\tilde u(t_\gamma), \gamma\tilde u(t_\gamma)]$ in a point $\tilde v(t'_\gamma)$.
%	 
%	\begin{figure}[htbp]
%		\centering
%		\begin{normalsize} 
%			\def\svgwidth{\textwidth}
%			%\def\svgwidth{20cm}
%			%\def\svgscale{0.5}
%			\input{{figs/}dessin6.pdf_tex}
%		\end{normalsize}
%		\caption{Disposition of the rays $\tilde u(\R^+)$ and $\tilde v(\R^+)$.}
%		\label{dessin6}
%	\end{figure}
%	
%	
%	
%	We shall remark that, because $d(i,(\gamma^-,\gamma^+))\geq 1$, we have $d(i,\tilde u(t_\gamma))\geq 1$. In addition, we showed that $\tilde v(\mathbb{R}^+)\cap [\tilde u(t_\gamma), \gamma\tilde u(t_\gamma)]\supset\{\tilde v(t'_\gamma)\}$, so $d(\tilde u(t_\gamma), \tilde v(t'_\gamma))\leq \ell(\gamma)$, which means we can choose the length $ \ell (\gamma)$ to be arbitrarily small so that $\tilde v(\mathbb{R}^+)$ is arbitrarily near $\tilde u(\mathbb{R}^+)$. In other words, the smaller $\ell(\gamma)$, the bigger $\tilde v(+\infty)$. 
%	
%	In order to study $\Delta_t$, we will first bound:
%	$$\Delta_t^d:=\min\{d(g_{t+r}\tilde v, g_t\tilde u), d(g_{t+r}\tilde v, \gamma g_t\tilde u)\}<\eta,
%	$$
%	and then we will study the missing terms in order to get the bound for $\Delta_t$. To do so, we will distinguish two separate cases:
%	\begin{itemize}
%		
%	
%		
%		\item \underline{Case 1}: %(\textcolor{blue}{Aqui referenciar la figura explicativa de este caso, aunque puede que no sea muy relevante en realidad}) 
%		$t\in [0,t'_\gamma]$
%		
%		We will show that, if $\ell(\gamma)$ is small enough, then $d(\tilde v(t+r), \tilde u(t))<\eta$. 
%		
%		We want to bound the following decomposition:
%		\begin{equation}\label{eq_caso1}
%			d(\tilde v(t+r),\tilde u(t))\leq d(\tilde v(t),\tilde v(t+r))+d(\tilde u(t),\tilde v(t)).
%		\end{equation}
%		Observe that, by aplying Proposition \ref{Corollary_1.5.5} we have that $r\leq \ell(\gamma)$. So,
%		
%		\begin{equation} \label{eq1}
%			d(\tilde v (t), \tilde v (t+r))=|r|\leq \ell (\gamma),
%		\end{equation}
%		%which concludes the bound for the first term. 
%		
%		
%		For the second term we will first observe that, by applying the triangle inequality to the points $i,\tilde u(t_\gamma)$ and $\tilde v(t'_\gamma)$ in Figure \ref{dessin6}, we obtain the following inequality:
%		\begin{equation}\label{eq2}
%			|t_\gamma-t'_\gamma|\leq \ell(\gamma).
%		\end{equation}
%		Since $t<t'_\gamma,$ $\tilde u(0)=\tilde v(0)=i$, and $\tilde v(+\infty)\neq \tilde u (\infty)$, it follows from the convexity of the function $d(\tilde u(\cdot),\tilde v(\cdot))$ that 
%		$d(\tilde u(t),\tilde v(t))\leq d(\tilde u(t'_\gamma),\tilde v(t'_\gamma))$, thus:
%		\begin{align*}\label{eq3}
%			d(\tilde u(t),\tilde v(t))\leq& d(\tilde u(t'_\gamma),\tilde v(t'_\gamma))\\
%			\leq &d(\tilde u(t'_\gamma),\tilde u(t_\gamma))+d(\tilde u(t_\gamma),\tilde v(t'\gamma))\\
%			\leq &2\ell(\gamma), 
%		\end{align*}
%		where the third inequality follows from inequality \ref{eq2} and the fact that $\tilde v(t'_\gamma)$ belongs to the hyperbolic segment $[\tilde u(t_\gamma),\gamma\tilde u(t_\gamma)]$.
%		
%		Finally, we conclude by applying these bounds in decomposition \ref{eq_caso1}:
%		
%		\begin{align*}
%		d(\tilde v(t+r),\tilde u(t)) &\leq d(\tilde v(t),\tilde v(t+r))+d(\tilde u(t),\tilde v(t))\\
%		& \leq \ell(\gamma)+2\ell(\gamma)=3\ell(\gamma).
%		\end{align*}
%		
%		\item \underline{Case 2}: $t\geq t'_\gamma$
%		
%		We will show that if $\ell(\gamma)$ is small enough, then $d(\tilde v(t+r), \gamma \tilde u(t))<\eta$. Let us define $s_t$ the time such that $d(\tilde v(t), \gamma \tilde u (s_t))=d(\tilde v(t),\gamma \tilde u(\mathbb R))$.  We will now distinguish two cases:
%		
%		\begin{itemize}
%			\item \underline{Case 2 (a)}: 
%			$t_\gamma\leq s_t$ 
%			
%			We have the following decomposition: 
%			
%			\begin{equation}\label{eq5}
%				d(\tilde v(t+r), \gamma \tilde u(t))\leq d(\tilde v(t+r), \tilde v(t))+d(\tilde v(t), \gamma \tilde u (s_t)) + d(\gamma\tilde u(s_t), \gamma \tilde u (t)).
%			\end{equation}
%			Let us bound each term. 
%			
%			The first term can be limited as before by applying Proposition  \ref{Corollary_1.5.5}:
%			
%			\begin{equation}\label{eq6}
%				d(\tilde v(t+r), \tilde v(t))=|r| \leq \ell(\gamma).
%			\end{equation}	
%			For the second term, since $\tilde v(+\infty)=\gamma \tilde u(+\infty)$ and $t\geq t_\gamma'$, we have
%			
%			\begin{equation}\label{eq7}
%				d(\tilde v(t), \gamma\tilde u(s_t))\leq d(\tilde v(t_\gamma '), \gamma \tilde u(\mathbb R))\leq d(\tilde v(t_\gamma '), \gamma \tilde u(t_\gamma))\leq \ell(\gamma), 
%			\end{equation}
%			
%			
%			\begin{figure}[htbp]
%				\centering
%				\begin{normalsize} 
%					\def\svgwidth{\textwidth}
%					%\def\svgwidth{20cm}
%					%\def\svgscale{0.5}
%					\input{{figs/}dessin17.pdf_tex}
%				\end{normalsize}
%				\caption{Up to conjugation, the rays $\tilde u(\R^+)$ and $\tilde \gamma^{-1}v(\R^+)$ are disposed this way.}
%				\label{dessin17}
%			\end{figure}
%
%			
%			The first inequality follows from the fact that there exists a $\delta-$neighborhood of $\tilde u(\R)$ that contains $\gamma^{-1}\tilde v(t)$ but not $\gamma^{-1}\tilde v(t'_\gamma)$ as we can see in Figure \ref{dessin17} and remark that $d(\tilde v(t'_\gamma),\gamma \tilde u(\R^+)= d(\gamma^{-1}\tilde v(t'_\gamma), \tilde u(\R^+) )$.
%	
%			Let us now limit the last term of \ref{eq5}. 
%			
%			$d(\gamma \tilde u(s_t),\gamma \tilde u(t))=|t-s_t|.$
%			Let $t'$ be the time such that $t=t'_\gamma+t'$ and $s_t'$ such that $s_t=t_\gamma+s_t'$. thanks to the assumption $s_t\geq t_\gamma$, we know that $s_{t}'\geq 0$ and that we are indeed in the situation described in Figure \ref{dessin7}. By applying the triangular inequality in the portion of figure \ref{dessin7} limited by the points $\gamma \tilde u(t_\gamma), \gamma\tilde u(s_t),\tilde v(t)$ and $\tilde v(t'_\gamma)$.
%			
%			$$t'-2\ell(\gamma)\leq s_t'\leq \ell (\gamma)+t'+\ell (\gamma)=t'+  2\ell (\gamma).$$
%			
%			By applying triangular inequality in the hyperbolic triangle defined by the points $i,\tilde u(t_\gamma)$ and $\tilde v(t'_\gamma)$ (see Figure \ref{dessin7}), we now obtain:
%			
%			\begin{figure}[htbp]
%				\centering
%				\begin{normalsize} 
%					\def\svgwidth{\textwidth}
%					%\def\svgwidth{20cm}
%					%\def\svgscale{0.5}
%					\input{{figs/}dessin7.pdf_tex}
%				\end{normalsize}
%				\caption{Majoration for the Case 2 (a).}
%				\label{dessin7}
%			\end{figure}
%			
%			
%			\begin{equation}\label{eq8}
%				t'_\gamma-\ell (\gamma)\leq t_\gamma\leq t_\gamma '+\ell (\gamma).
%			\end{equation}
%			
%			Finally, by adding both inequalities we have
%			$$t'+t_\gamma'-3\ell(\gamma)= t-3\ell(\gamma)\leq s_t\leq t'+t_\gamma'+3\ell(\gamma)=t+3\ell(\gamma),$$
%			from which we deduce:
%			
%			$$|s_t-t|\leq 3\ell(\gamma).$$ 
%			And we conclude Case 2 (a) by applying these bounds in inequality \ref{eq5}:
%			\begin{align*}
%			d(\tilde v(t+r), \gamma \tilde u(t)) &\leq d(\tilde v(t+r), \tilde v(t))+d(\tilde v(t), \gamma \tilde u (s_t)) + d(\gamma\tilde u(s_t), \gamma \tilde u (t)) \\
%			&\leq \ell(\gamma) +\ell(\gamma) +3\ell(\gamma)=5\ell(\gamma)
%			\end{align*}
%			 
%			
%			
%			\item \underline{Case 2 (b)} $t_\gamma> s_t$
%			
%			In order to study this case we first define $T\equiv T(\gamma)$ as the time such that $d(\tilde v(T),\gamma \tilde u (\R^))=d(\tilde v(T), \gamma \tilde u(t_\gamma))$, i.e, $s_T=t_\gamma$. Observe that  $t_\gamma> s_t$ is equivalent to $t'_\gamma\leq t< T$. (See Figure \ref{dessin16})
%			
%			\begin{figure}[htbp]
%				\centering
%				\begin{normalsize} 
%					\def\svgwidth{\textwidth}
%					%\def\svgwidth{20cm}
%					%\def\svgscale{0.5}
%					\input{{figs/}dessin16.pdf_tex}
%				\end{normalsize}
%				\caption{Majoration for the Case 2 (b).}
%				\label{dessin16}
%			\end{figure}
%			
%			
%			We will consider the following decomposition: 
%			\begin{align}\label{eq_descom}
%				d(\gamma \tilde u(t),\tilde v(t+r))& \leq d(\gamma \tilde u(t),\tilde v(t))+d(\tilde v(t),\tilde v(t+r)) \nonumber \\
%				& \leq d(\gamma\tilde u(t),\tilde v(t))+ \ell(\gamma).
%			\end{align}
%			Therefore, our aime is to find a superior bound for:
%			\begin{align}\label{eq18}
%				d(\gamma \tilde u(t),\tilde v(t))& \leq \,d(\gamma \tilde u(t),\gamma\tilde u(T))+d(\gamma\tilde u(T),\tilde v(T))+ d(\tilde v(T),\tilde v(t))\nonumber \\ & \leq \, 2|T-t|+ 5\ell(\gamma) \nonumber\\ 
%				& \leq  2|T-t'_\gamma|+ 5\ell(\gamma),
%			\end{align}
%			where inequality $d(\gamma\tilde u(T),\tilde v(T))\leq 5\ell (\gamma)$ follows from Case 2 (a) as $s_T=t_\gamma$. We then need to majorate $|T-t'_\gamma|$.
%			
%			To do so, let us observe that there exists a $\delta-$neighborhood of $\tilde u(\R)$ that contains $\gamma^{-1}\tilde v(t'_\gamma)$ as we can see in Figure\ref{dessin16}. Thus:
%			
%			\begin{align}\label{eq4}
%				d(\gamma^{-1}\tilde v(T),\tilde u(t_\gamma))=&d(\gamma^{-1}\tilde v(T),\tilde u(\R))\leq d(\gamma^{-1}\tilde v(t'_\gamma),\tilde u(\R))  \nonumber \\
%				\leq &d(\gamma^{-1}\tilde v(t'\gamma),\tilde u(t_\gamma)\leq \ell(\gamma)).
%			\end{align} 
%			And then, by applying the triangular inequality in the hyperbolic triangle defined by the points $\gamma^{-1}\tilde v(t'_\gamma)$, $\gamma^{-1}\tilde v(T)$ and $\tilde u(t_\gamma)$  (see Figure \ref{dessin16})
%			\begin{align*}
%				|T-t'_\gamma|& =  d(\gamma^{-1}\tilde v(t'_\gamma),\gamma^{-1}\tilde v(T)) \\ \, & \leq d(\gamma^{-1}\tilde v(t'_\gamma),\tilde u(t'_\gamma)+d(\tilde u(t'_\gamma),\gamma^{-1}\tilde v(T))\\
%				& \leq 2\ell(\gamma),
%			\end{align*}
%			
%			where $d(\tilde u(t'_\gamma),\gamma^{-1}\tilde v(T))\leq \ell(\gamma)$ follows from inequality \ref{eq4}. Now, replacing in inequality \ref{eq18} we obtain 
%			
%			$$d(\gamma \tilde u(t),\tilde v(t))\leq 9\ell(\gamma).$$ 
%			And we conclude the Case 2 (b) by applying this bound in decomposition \ref{eq_descom}:
%			\begin{align*}
%				d(\gamma \tilde u(t),\tilde v(t+r))& \leq d(\gamma \tilde u(t),\tilde v(t))+d(\tilde v(t),\tilde v(t+r)) \nonumber \\
%				& \leq d(\gamma \tilde u(t),\tilde v(t))+ \ell(\gamma)\\
%				& \leq  9\ell(\gamma)+\ell(\gamma)=10\ell(\gamma).
%			\end{align*}
%			
%			
%			
%			
%		\end{itemize}
%		This concludes the bounds for $\Delta_t^d$.
%	\end{itemize}
%	
%
%	Let us now compute the bounds for $\Delta_t$ by studying the distance $d_1$. 
%	We can distinguish three scenarios:
%	\begin{enumerate}[(a)]
%		\item If $t+1<t_\gamma '$, then we are in first case and we have that both $d(\tilde v(t+r+1), \tilde u(t+1))$ and $d(\tilde v(t+r), \tilde u(t))$ are bouded by $3\ell(\gamma)$. Then $d_1(g_{t+r}\tilde v, g_t\tilde u\leq 6 \ell(\gamma)$.
%		
%		\item If $t\geq t'_\gamma$ then so is $t+1$, and we are in the Case 2, so we can still bound both $d(\tilde v(t+r+1),\gamma \tilde u(t+1))$ and $d(\tilde v(t+r),\gamma \tilde u(t))$ by $10\ell(\gamma)$. Then $d_1(g_{t+r}\tilde v, g_t\tilde u\leq 20 \ell(\gamma)$.
%		
%		\item If $t+1\geq t'_\gamma$ and $t\leq t'_\gamma$, we will see that we find an upper bound for $d_1(g_{t+r}\tilde v, \gamma g_t\tilde u)$. Let us first notice that 
%		
%		\begin{align*}
%			d_1((g_{t+r}\tilde v, \gamma g_t\tilde u)) &\leq d(\tilde v(t+r), \gamma \tilde u(t)) + 
%			d(\tilde v(t+1+r), \gamma \tilde u(t+1))\\ 
%			&\leq  d(\tilde v(t+r), \gamma \tilde u(t))+5\ell(\gamma),
%		\end{align*}
%
%		where the second inequality from Case 2.
%		Thus, we have to bound \\ $d(\tilde v(t+r), \gamma \tilde u(t))$. Consider the following decomposition: 
%		
%		\begin{align}\label{eq9}
%			d(\tilde v(t+r), \gamma \tilde u(t)) &\leq  d(\tilde v(t+r), \tilde u(t)) + d(\tilde u(t), \gamma \tilde u(t)) \notag\\ 
%			& \leq 3\ell(\gamma) + d(\tilde u(t), \gamma \tilde u(t)), \notag
%		\end{align}
%		where the second inequality follows from Case 1. Therefore, we have to bound $d(\tilde u(t), \gamma \tilde u(t))$. 
%		We apply the formula \ref{formula_distancia} to $z=\tilde u(t)$: 
%		\begin{equation}\label{eq10}
%			\sinh\left(\frac{d(\tilde u(t),\gamma \tilde u(t))}{2}\right)=\cosh(d(\tilde u(t),(\gamma^-,\gamma^+)))\sinh\left(\frac{\ell(\gamma)}{2}\right).
%		\end{equation}
%		Now, since $t<t'_\gamma$ and $t+1>t'_\gamma$, we deduce that 
%		\begin{equation}\label{eq11}
%			0<|t-t'_\gamma|<1,
%		\end{equation}
%		Let us now see what happens to $d(\tilde u(t),(\gamma^-,\gamma+))$ in \ref{eq10}:
%		\begin{align*}
%			d(\tilde u(t),(\gamma^-,\gamma^+))&=|t-t_\gamma|=|t-t_\gamma+t_\gamma'-t_\gamma'|\\&\leq |t-t_\gamma'|+|t_\gamma'-t_\gamma|\leq 1+\ell(\gamma),
%		\end{align*}
%		%where the las inequality follows from equations \ref{eq11} and \ref{eq8}.
%		 Therefore equation \ref{eq10} is reduced to: 
%		\begin{align*}
%			\sinh\left(\frac{d(\tilde u(t),\gamma \tilde u(t))}{2}\right)
%			& \leq \cosh(1+\ell(\gamma))\sinh\left(\frac{\ell(\gamma)}{2}\right).
%		\end{align*}
%		
%		And because $\sinh$ is an increasing function that tends to $0$ as $\ell(\gamma)$ tend to $0$ we conclude the desired majoration.
%		
%%		Consider the function $\Phi$ defined as
%%		$$x\mapsto \Phi(x)= \operatorname{arcsinh}\left(2\cosh(1+x)\sinh\left(\frac{x}{2}\right)\right).$$
%%		This is a positive function that tends to $0$ as $x$ tends to zero. Thus, we have the following bound: 
%%		$$d(\tilde u(t),\gamma \tilde u(t)) \leq \Phi(\ell(\gamma)).$$
%		
%		%And we conclude by applying this bound in equation \ref{eq9}.
%		
%		
%		%Esto era para la flauta F1 pero ya lo podemos hacer para cualquier superficie hip�rbolica
%		%		To do so we will suppose two scenarios 
%		%		\begin{enumerate}[(i)]
%			%		\item If $t\geq t_\gamma$, then we can decompose $d(\tilde u(t), \gamma \tilde u(t))$ as follows: 			
%			%			\begin{align*}
%				%				d(\tilde u(t), \gamma \tilde u(t)) &\leq  d(\tilde u(t), \tilde u(t_\gamma)) + 	d(\tilde 	u(t_\gamma), 	\gamma \tilde u(t_\gamma)) + d(\gamma \tilde u(t_\gamma) \gamma \tilde u(t)), \\
%				%				& \leq |t_\gamma-t'_\gamma| + \ell(\gamma)+ 	|t_\gamma-t'_\gamma|, \\
%				%				& \leq 3\ell(\gamma),
%				%			\end{align*}
%			%			
%			%			where the first inequality follow from the fact that $\tilde u(t_\gamma)$ belongs to the axis $(\gamma^-,\gamma^+)$, and the second inequality follows from equation \ref{eq8}. 
%			%			
%			%		\item If $t< t_\gamma$. It is noteworthy that the function $t \mapsto d(\tilde u(t), \gamma \tilde u(t))	$ is decreasing in the interval $[0,t_\gamma]$ \textcolor{blue}{Para justificar esto podemos hacer referencia a la proposicion donde se da la formula de la distancia con respecto a la longitud de la isometria y la distancia al eje}. Hence: 
%			%		$$
%			%		d(\tilde u(t), \gamma \tilde u(t))\leq d(\tilde u(0), \gamma \tilde u(0)) = d(i, \gamma i).
%			%		$$ 
%			%		\end{enumerate}			
%		
%		
%	
%	\end{enumerate}
%	
%\end{proof}	

\subsection{Construction of suitable $\alpha=(\alpha_n)_{n\in\N }$}



\begin{proposition}\label{lema_clave}
	Let $\alpha=(\alpha_n=\gamma_{p_n})_{n\geq0}$ be a subsequence of $(\gamma_n)_{n\geq0}$ verifying $\ell({\alpha_n})\leq \frac{1}{2^{2n+3}}$. Then, 
	for every $n,\ell\geq 0$ there exists $T_\ell\ge 0$ with the property that,
	\[\forall \,t\ge T_\ell,~
	d_1\bigl(g_t g_{r_n}v_n,\, g_t u_0\bigr)
	\le \left(\sum_{k=0}^{n}\frac{1}{2^{k}}\right)\frac{1}{2^{\ell}} .
	\]
\end{proposition}

%It follows from this theorem that
%\[
%\forall \ell\geq 0,\ \exists T_\ell\ge 0\ \text{ such that }\ \forall t\ge T_\ell,\qquad
%d_1\bigl(g_t\omega_\alpha,\, g_tu_0\bigr)\le \frac{1}{2^{\ell-1}},
%\]
%and hence $\omega_\alpha\in W^{ss}(u_0)$.
%Applying Property \ref{property1}, one obtains:
%
%\medskip
%\noindent\textbf{Corollary.}
%\[
%\omega_\alpha\in W^{ss}(u_0)\setminus h_\R (u_0).
%\]

\subsubsection*{Proof of Proposition~\ref{lema_clave}}
For $m\ge 0$, we say that $\alpha_0,\dots,\alpha_m$ satisfy Property $(P_m)$ if:
\begin{itemize}
	%\item for every $n\in\{0,\dots,m\}$, $\ \ell(\alpha_n)\le \dfrac{1}{2^{n}}$;
	\item for every $n,\ell\in\{0,\dots,m\}$, there exists $T_l\ge 0$ such that for every $t\ge T_l$,
	\[ \forall t\ge T_l,\quad
	d_1\bigl(g_t g_{r_n} v_n,\, g_t u_0\bigr)
	\le \left(\sum_{k=0}^{n}\frac{1}{2^{k}}\right)\frac{1}{2^{\ell}} .
	\]
\end{itemize}

\medskip
\noindent\textbf{Step $(P_0)$.}
By definition, $\tilde v_0=\operatorname{Wind}_{\alpha_0}(\tilde u_0)$ and $r_0=\tau_{\alpha_0,\tilde v_0}$.
Recall that $\alpha_0=\gamma_{p_0}$ is such that $\ell(\alpha_0)\le \frac{1}{2^3}$.
%\begin{itemize}
%	%\item $d\bigl(i,(\alpha_0^-,\alpha_0^+)\bigr)\ge 1,$
%	\item $\tilde u_0(\mathbb R^+)\cap (\alpha_0^-,\alpha_0^+)\neq\varnothing,$
%	\item $\ell(\alpha_0)\le \frac{1}{2^3}.$
%\end{itemize}
Applying Key Proposition \ref{key_prop} one obtains $(P_0)$:
\begin{itemize}
%	\item $\ell(\alpha_0)\le 1$;
	\item$\forall t\ge T_0=0$,
	\[
	d_1\bigl(g_tg_{r_0} v_0,\, g_t u_0\bigr)\le 1.
	\]
\end{itemize}

\medskip
\noindent\textbf{Induction step.}
Let $m\geq0$, and suppose that $(P_m)$ is satisfied.
Since $|r_0|\le \ell(\alpha_0)\le 1$ (Proposition \ref{bound_wind}) and
$|r_n-r_{n-1}|\le \ell(\alpha_n)\le \frac{1}{2^{n}}$ (Proof of Proposition \ref{lema_4}), we have
\[
\forall n\in\{0,\dots,m\},\qquad
|r_n|\le 1+\sum_{k=0}^{n}\frac{1}{2^{k}}\le 3.
\]

\medskip
Recall that $\alpha_{m+1}=\gamma_{p_{m+1}}$ is such that $\ell(\alpha_{m+1})\le \frac{1}{2^{2(m+1)+3}}$. We have to prove the existence of	
$T_{m+1}\ge 0$ satisfying:
\begin{enumerate}
	\item for every $n\in\{0,\dots,m\},$ 
	\[\forall \, ~ t\ge T_{m+1},~
	d_1\bigl(g_t g_{r_{n}}v_n,\, g_tu_0\bigr)
	\le \left(\sum_{k=0}^{m}\frac{1}{2^{k}}\right)\frac{1}{2^{m+1}};
	\]
	\item for every $\ell\in\{0,\dots,m+1\},$ 
	\[\forall \, ~ t\ge T_{\ell},~
	d_1\bigl(g_t g_{r_{m+1}}v_{m+1},\, g_tu_0\bigr)
	\le \left(\sum_{k=0}^{m+1}\frac{1}{2^{k}}\right)\frac{1}{2^{\ell}} .
	\]
\end{enumerate}

\medskip
\noindent\textbf{Case (1).}
By construction, there exists $s_n\in \R$ such that $\beta_n^{-1}g_{r_n}\tilde v_n= h_{s_n}(\tilde u_0)$.
Since $g_th_sg_{-t}=h_{se^{-t}}$, we have that
$$
d_1(g_th_{s_n}\tilde u_0,g_t\tilde u_0)=d_1(h_{s_ne^{-t}}g_t\tilde u_0,g_t\tilde u_0).
$$
Since $d((h_{s_ne^{-t}}g_t\tilde u_0)(0),\tilde u(t))\leq |s_n|e^{-t}$, we obtain
$$d_1(g_th_{s_n}\tilde u_0 ,g_t \tilde u_0)\leq |s_n|(e^{-t}+e^{-(t+1)}).$$

Let $t_n\geq 0$ such that
$$
\forall\, t\geq t_n, ~d_1(g_th_{s_n}\tilde u_0,g_t\tilde u_0)\leq \frac{1}{2^{m+1}}.
$$
Take $T_{m+1}'=\max\{t_0,\dots,t_m\}$: 	
\[\forall \, t\ge T_{m+1}', ~
d_1\bigl(g_t\beta_n^{-1}g_{r_n}\tilde v_n,\, g_t\tilde u_0\bigr)
\le \frac{1}{2^{m+1}}.
\]
This implies
\begin{equation}\label{case1}
	\forall \, t\ge T_{m+1}', ~
	d_1\bigl(g_t g_{r_n} v_n,\, g_t u_0\bigr)
	\le \left(\sum_{k=0}^{m}\frac{1}{2^{k}}\right)\frac{1}{2^{m+1}}.
\end{equation}

\medskip
\noindent\textbf{Case (2).}
We have
\[
d_1\bigl(g_t g_{r_{m+1}} v_{m+1},\, g_t u_0\bigr)
\le
\underbrace{d_1\bigl(g_t g_{r_{m}} v_{m},\, g_t u_0\bigr)}_{(a)}
+
\underbrace{d_1\bigl(g_t g_{r_{m+1}} v_{m+1},\, g_t g_{r_{m}} v_{m}\bigr)}_{(b)}.
\]

\medskip
\noindent\textbf{Part (a).}
Since $(P_m)$ is satisfied, for any $\ell\in\{0,\dots,m\}$ there exists $T_\ell\ge 0$ such that for every $t\ge T_\ell$,
\[
(a)\le \left(\sum_{k=0}^{m}\frac{1}{2^{k}}\right)\frac{1}{2^{\ell}} .
\]
Moreover, by (\ref{case1}), we have
\[ \forall\, t\ge T_{m+1}^{'}:~
(a)\le \left(\sum_{k=0}^{m}\frac{1}{2^{k}}\right)\frac{1}{2^{m+1}}.
\]

\medskip
\noindent\textbf{Part (b).}
Recall that $\tilde v_{m+1}=\operatorname{Wind}_{\beta_m \alpha_{m+1}\beta_m^{-1}}(\tilde v_m)$ and
\[
\tau_{\beta_m \alpha_{m+1}\beta_m^{-1},\,\tilde v_m}=r_{m+1}-r_m.
\]
Since $\ell(\alpha_{m+1})\leq \frac{1}{2^{2(m+1)+3}}$, applying Key Proposition \ref{key_prop}, we obtain:
\[\forall \, t\geq 0,~
d_1\bigl(g_{t+r_{m+1}-r_m}\tilde v_{m+1},\, g_t\tilde v_{m}\bigr)
\le \frac{1}{2^{2(m+1)}}.
\]
Hence, for every $s>-r_m$,
\[
d_1\bigl(g_s g_{r_{m+1}} v_{m+1},\, g_s g_{r_m} v_m\bigr)
\le \frac{1}{2^{2(m+1)}}.
\]
Since $|r_m|\le 3$, one has for every $t\ge 3$,
\[
d_1\bigl(g_t g_{r_{m+1}} v_{m+1},\, g_t g_{r_m} v_m\bigr)
\le \frac{1}{2^{2(m+1)}}.
\]

\medskip
Finally, let
\[
T_{m+1}=\max\{3,\,T_{m+1}'\}.
\]
Since for every $\ell\in\{0,\dots,m+1\}$,
\[
\frac{1}{2^{2(m+1)}}\le \frac{1}{2^{m+1}}\cdot\frac{1}{2^{\ell}},
\]
one obtains for every $\ell\in\{0,\dots,m+1\}$ and every $t\ge T_\ell$,
\[
d_1\bigl(g_t g_{r_{m+1}} v_{m+1},\, g_t u_0\bigr)
\le \left(\sum_{k=0}^{m+1}\frac{1}{2^{k}}\right)\frac{1}{2^{\ell}}.
\]
\qed


\begin{remark}
Since the only condition imposed on $\alpha_{m+1}$ at the Induction Step is \\ $\ell(\alpha_{m+1})\leq \frac{1}{2^{2(m+1)+3}}$ and since the sequence $(\ell(\gamma_n))_{n\geq 0}$ is decreasing  and converges to $0$, it follows that the set of such sequences $\Sigma$ is uncountable.  

\end{remark}





\subsection{End of the Proof of the Main Theorem}\label{end_proof}


\begin{proof}
		We denote $\Sigma$ the set of $(\alpha_n)_{n\geq 0}$ with $\ell(\alpha_n)\leq \frac{1}{2^{(2n+3)}}$.
		Let $(\alpha_n)_{n\in\N}$ a subsequence in $\Sigma$ (see Proposition \ref{lema_clave}). Referring back to the notations of Lemma \ref{lemma_vectors} and Proposition \ref{lema_4}, we will prove that the vector $w_\alpha=g_{r_\alpha} v_\alpha$ of $T^1S_0$ belongs to $W^{ss}u_0.$ 
		Observe that, thanks to Proposition \ref{lema_clave}, for all $n,l\geq 0$ and $t\geq T_l:$
		$$ d_1(g_{t+r_n}v_n,g_t u)<\frac{1}{2^{l-1}}$$ And thus, taking the limit on $n$, we can state that for all $t\geq T_l$:
		$$ d_1(g_{t+r}v_\alpha,g_t u_0)<\frac{1}{2^{l-1}}, $$
		then, for all $\varepsilon >0 $ there exists $l\geq 0$ such that for all $t\geq T_l$:
		$$ d_1(g_tw_\alpha,g_tu_0)=d_1(g_{t+r}v_\alpha,g_tu_0)\leq\varepsilon,$$
		which implies that $w\in W^{ss}u$.

Let
$
W(\infty)=\left\{\tilde w_{\alpha}(+\infty),\ \alpha\in\Sigma\right\}.
$
This set is uncountable. Suppose it is not, then there exists a sequence
$
\bigl(\tilde w_{\alpha^i}(+\infty)=x_i\bigr)_{i\in\mathbb N}
$
with $\alpha^i\in\Sigma$ such that
$
W(\infty)=\{x_i,\ i\in\mathbb N\}.
$
Let us construct $\alpha'\in\Sigma$ such that
$
\tilde w_{\alpha'}(+\infty)\neq x_i
$
for any $i\in\mathbb N$.

Choose
$
\alpha_0'=\gamma_{p_{k_0}}
$
such that:
\begin{itemize}
	\item $\ell(\alpha_0')\leq \dfrac{1}{2^3}$,
	\item $
	\alpha_0^{\prime +}>x_0.
$
\end{itemize}

By induction, if $\alpha_0',\ldots,\alpha_n'$ are chosen, take
$
\alpha_{n+1}'
$
satisfying:
\begin{itemize}
	\item $\ell(\alpha_{n+1}')\leq \dfrac{1}{2^{2(n+1)+3}}$,
	\item $
	\alpha_{n+1}^{\prime +}>
	\alpha_n^{\prime -1}\cdots \alpha_0^{\prime -1}(x_{n+1}).
	$
\end{itemize}

This construction is possible since
$
\lim\limits_{n\to+\infty}\gamma_n^+=\infty
$
and the sequence
$
\bigl(\ell(\gamma_n)\bigr)_{n\geq 0}
$
decreases to $0$.

Observe that, using the position of the axis
$
(\gamma_n^-,\gamma_n^+) \text{ for } n\geq 0
$
and the dynamics of $\gamma_n$, for any $0\leq n\leq m$, we have:
\[
\alpha_n'\ \alpha_{n+1}'\ \ldots\ \alpha_m'(\infty)>\alpha_n^{\prime +}.
\]

We deduce for $n\geq 1$:
\[
\alpha_{n-1}^{\prime -1}\cdots \alpha_0^{\prime -1}
\bigl(\tilde w_{\alpha'}(+\infty)\bigr)
>
\alpha_n^{\prime +}
\qquad (*)
\]

By construction, for $n\geq 1$, we have
\[
\alpha_n^{\prime +}>
\alpha_{n-1}^{\prime -1}\cdots \alpha_0^{\prime -1}(x_n)
\qquad (**)
\]

Both inequalities imply:
\[
\forall n\geq 1,\qquad
\alpha_{n-1}^{\prime -1}\cdots \alpha_0^{\prime -1}
\bigl(\tilde w_{\alpha'}(\infty)\bigr)
\neq
\alpha_{n-1}^{\prime -1}\cdots \alpha_0^{\prime -1}(x_n),
\]
and hence
$
\tilde w_{\alpha'}(+\infty)\neq x_n.
$
Moreover
$
\tilde w_{\alpha'}(+\infty)\neq x_0
$
since
$
\alpha_0^{\prime +}>x_0
\text{ and }
\tilde w_{\alpha'}(+\infty)>\alpha_0^{\prime +}.
$

We conclude that $W(\infty)$ is not countable.

In particular, since $\Gamma$ is countable, the set
$
\Gamma W(\infty)=\{\gamma(x),\ \gamma\in\Gamma,\ x\in W(\infty)\}
$
is an uncountable disjoint union of $\Gamma$-orbits.

Let $v\in W^{ss}(u_0)$, since $v\in h_\R(u_0)$
implies
$
\tilde v(\infty)\in \Gamma\infty,
$
the set $W^{ss}(v_0)$ is an uncountable union of horocycle trajectories.


\end{proof}


	
	
%\appendix 
%\section{Equivalence between different distances in the unitary tangent bundle}
%\label{apendice}
%%todas las notaciones deber�amos definirlas en unos preliminares
%
%%	We denote distance of the hyperbolic plane by $d$. For $v\in T^1 \mathbb{H}$, $c_v$ is the unit speed geodesic with $c_v'(0)=v$. 
%Let $\tilde v,\tilde w\in T^1 \mathbb{H}^2$. We recall %ya definida
%$$d_1(\tilde v,\tilde w):=d(\tilde v(0),\tilde w(0))+d(\tilde v(1),\tilde w(1)).$$
%
%Let $c_{\tilde v,\tilde w}$ be the geodesic passing through the base points of $\tilde v$ and $\tilde w$. Let $\alpha, \beta\in [-\pi,\pi)$ be the oriented angles formed by $\tilde v$ and $\tilde w$ with the geodesic $c_{\tilde v,\tilde w}$, respectively.
%We define $$d_2(\tilde v,\tilde w):=d(\tilde v(0),\tilde w(0))+|\alpha-\beta|,$$
%
%which is equivalent to the Sasaki metric defined as: 
%$$
%d_S(\tilde v,\tilde w)= \sqrt{d(\tilde v(0),\tilde w(0))^2+|\alpha-\beta|^2}.
%$$
%
%\begin{proposition}\label{quadrilateral}
%	Let $\tilde v,\tilde w\in T^1 \mathbb{H}$. We set $D_0:=d(\tilde v(0),\tilde w(0))$, $D_1:=d(\tilde v(1),\tilde w(1))$, $S:=\sinh(1)$ and $C:=\cosh(1)$. They satisfy
%	\begin{align}\label{quadrilateralequation}
%		\begin{split}
%			\cosh D_1=&C^ 2 \cosh D_0 -2CS\sinh D_0 \sin(\frac{\alpha-\beta}{2}) \sin(\frac{\alpha+\beta}{2})
%			\\
%			&-S^ 2 \cosh^ 2(\frac{D_0}{2})\cos(\alpha-\beta)-S^ 2\sinh^ 2(\frac{D_0}{2})\cos(\alpha+\beta).
%		\end{split}
%	\end{align}
%\end{proposition}
%\begin{proof}
%	We consider the triangle $T_1$ with vertices $\tilde v(0),\tilde w(0)$ and $\tilde v(1)$ and the triangle $T_2$ with vertices $\tilde w(0)\tilde ,v(0$ and $\tilde w(1)$. Let $\beta'$ the interior oriented angle of $T_1$ at $\tilde w(0)$. $T_1$ has sides of length $D_0$, $1$ and $D_2:=d(\tilde v(1), \tilde w(0))$. Applying the hyperbolic cosine rule for $T_1$, we have 
%	\begin{equation}\label{eq:1}
%		\cosh D_2= C \cosh D_0 -S \sinh D_0 \cos(\pi - \alpha),
%	\end{equation}
%	\begin{equation}\label{eq:2}
%		C=\cosh D_0 \cosh D_2 -\sinh D_0 \sinh D_2 \cos \beta'.
%	\end{equation}
%	By the sine rule between the angles at points $v(0)$ and $w(0)$, we also have
%	\begin{equation}\label{eq:3}
%		\sinh D_2 \sin \beta '= S \sin(\pi-  \alpha ).
%	\end{equation}
%	The angle of $T_2$ at $w(0)$ is $\beta-\beta'$. By the hyperbolic cosine rule for $T_2$, 
%	\begin{equation}\label{eq:4}
%		\cosh D_1 = C\cosh D_2 -S \sinh D_2 \cos(\beta-\beta').
%	\end{equation}
%	
%	Replacing (\ref{eq:1}) in (\ref{eq:2}), we obtain 
%	\begin{equation*}
%		C=C\cosh^ 2 D_0 +S\cosh D_0 \sinh D_0 \cos \alpha -\sinh D_0 \sinh D_2 \cos \beta ',
%	\end{equation*}
%	which can be rewritten as
%	\begin{equation}\label{eq:5}
%		\sinh D_2 \cos \beta'  = C\sinh D_0 +S\cosh D_0 \cos \alpha.
%	\end{equation}
%	
%	By a trigonometric relation at (\ref{eq:4}) and then replacing  (\ref{eq:1}), (\ref{eq:3}) and (\ref{eq:5}), we have 
%	\begin{align*}
%		\cosh D_1 &= C\cosh D_2 -S \sinh D_2 \cos \beta \cos\beta' -S \sinh D_2 \sin \beta \sin\beta' \\
%		&=C^ 2\cosh D_0 + CS \sinh D_0 \cos \alpha -CS\sinh D_0 \cos \beta \\
%		&-S^2 \cosh D_0 \cos \alpha \cos \beta - S^ 2 \sin \alpha \sin \beta. 
%	\end{align*}
%	Elementary trigonometric manipulations lead to the formula of the statement.
%\end{proof}
%
%
%\begin{lemma}\label{lemma}
%	There exists $\varepsilon>0, K>0$ such that, for all $\alpha, \beta \in \R $, $D_0,D_1\in \R_+$ satisfying Equation (\ref{quadrilateralequation}), if $D_0+|\alpha-\beta|<\varepsilon$ then
%	$$D_1\le K(D_0+|\alpha-\beta|)$$ 
%	and if $D_0+D_1<\varepsilon$ then
%	$$|\alpha-\beta|\le D_0+D_1.$$
%\end{lemma}
%
%\begin{proof}
%	We can compare the Taylor polynomials of order two at $D_1=D_0=\alpha-\beta=0$ of both sides of (\ref{quadrilateralequation}). We have
%	\begin{align*}
%		1+\frac{D_1^ 2}{2}+o_2(D_1)&= C^ 2(1+\frac{D_0^ 2}{2})-CS D_0 (\alpha-\beta) \sin(\frac{\alpha+\beta}{2}) \\
%		&- S^ 2(1+\frac{D_0^ 2}{8})^ 2 (1-\frac{(\alpha-\beta)^ 2}{2})-\frac{1}{4}S^ 2 D_0^ 2\cos (\alpha+\beta) +o_2(D_0,\alpha-\beta) \\
%		&=1 +\frac{1}{2}\left(C^ 2-S^ 2\frac{1+\cos(\alpha+\beta)}{2}\right) D_0^ 2 -CS\sin \left(\frac{\alpha+\beta}{2}\right) D_0 (\alpha-\beta) \\
%		&+\frac{1}{2}S^ 2(\alpha-\beta)^ 2 +o_2(D_0,\alpha-\beta).
%	\end{align*}
%	The polynomial of the right hand side is bounded below by $1+\frac{S^ 2}{2}(|\alpha-\beta |-D_0)^ 2$ and above  by $1+\frac{C^ 2}{2}(|\alpha-\beta |+D_0)^ 2$. Now we choose constants $K,K'$ such that $1<K'<S<C<K$. There exists $\varepsilon>0$ depending only on $K$ and $K'$ small enough so higher order terms can be neglected: if $D_0+|\alpha-\beta|<\varepsilon $ then 
%	\begin{align*}
%		1+\frac{D_1^ 2}{2}\le 1+\frac{K^ 2}{2}(|\alpha-\beta|+D_0)^ 2
%	\end{align*}
%	and if $D_0+D_1<\varepsilon$
%	\begin{align*}
%		1+\frac{K'^ 2}{2}(|\alpha-\beta|- D_0)^ 2\le 	1+\frac{D_1^ 2}{2}.
%	\end{align*}
%	The first inequality is equivalent to $D_1\le K(D_0+|\alpha-\beta|)$ and the second one is equivalent to $|\alpha -\beta|\le D_0+ \frac{1}{K'}D_1$, which implies the statement.
%\end{proof}
%
%\begin{proposition}
%	The distances $d_1$ and $d_2$ are equivalent. 
%\end{proposition}
%
%\begin{proof}
%	We have $D_1\le d(\tilde v(0),\tilde v(1))+D_0 +d(\tilde w(0),\tilde w(1))$.
%	If $D_0+|\alpha-\beta|\ge \varepsilon$, then $$
%	D_1\le D_0+\frac{2}{\varepsilon}(D_0+|\alpha-\beta|)�
%	$$
%	In view of Lemma \ref{lemma}, we have 
%	$$ d_1(\tilde v,\tilde w) \le D_0 + \max (D_0+\frac{2}{\varepsilon}(D_0+|\alpha-\beta|), K (D_0 +|\alpha-\beta|)) \le \max(2+\frac{2}{\varepsilon}, K+1)d_2(\tilde v,\tilde w).$$ 
%	
%	On the other hand, if $D_0+D_1\ge \varepsilon$, then $|\alpha-\beta|\le 2\pi \le \frac{2\pi}{\varepsilon}(D_0+D_1)$. Therefore
%	$$d_2(\tilde v,\tilde w)\le D_0 +\max (\frac{2\pi}{\varepsilon},1)(D_0+D_1)\le \max(1+\frac{2\pi}{\varepsilon}, 2)d_1(\tilde v,\tilde w).
%	$$
%\end{proof}
%
%\begin{remark}
%	Since the metrics $d_1,d_2$ and $d_S$ are equivalent,  the strong-stable set may be defined indifferently with either distance. Equivalently, for any $u\in T^1 S$:
%	$$
%	W^{ss}u=\{ v\in T^1S:~ \lim_{t\to \infty} d_S(g_tu,g_tv)=0 \}= \{ v\in T^1S:~ \lim_{t\to \infty} d_1(g_tu,g_tv)=0 \}.
%	$$
%\end{remark}

\appendix

\section{Equality between $W_{d_1}^{ss}(u)$ and $W_{d_2}^{ss}(u)$}\label{apendice}

Let $\tilde v,\tilde w\in T^{1}\mathbb H^{2}$ such that $\tilde v(0)\neq \tilde w(0)$.
Consider the geodesic ray starting at $\tilde v(0)$ and passing through $\tilde w(0)$.
Denote $C_{\tilde v,\tilde w}\in T^{1}\mathbb H^{2}$ its unitary tangent vector at $\tilde v(0)$.

By construction, for $s=d(\tilde v(0),\tilde w(0))$, the base point of $g_{s}C_{\tilde v,\tilde w}$ is $\tilde w(0)$.
Denote $\alpha_{\tilde v,\tilde w}\in[0,2\pi)$ (resp.\ $\beta_{\tilde v,\tilde w}$) the oriented angle
defined by $\widehat{(C_{\tilde v,\tilde w},\tilde v)}$ (resp.\ $\widehat{\left(g_{s}(C_{\tilde v,\tilde w}),\tilde w\right)}$).

We introduce a distance
$ 
d_{2}:T^{1}\mathbb H^{2}\times T^{1}\mathbb H^{2}\longrightarrow \mathbb R_{+}
$
defined by
\[
d_{2}(\tilde v,\tilde w)=d\bigl(\tilde v(0),\tilde w(0)\bigr)+\bigl|\alpha_{\tilde v,\tilde w}-\beta_{\tilde v,\tilde w}\bigr|,
\]
where $d$ is the hyperbolic distance.

It turns out that $d_{2}$ is equivalent to the Sasaki distance $d_{Sa}$ induced by the
Sasaki metric on $T^{1}\mathbb H^{2}$ \cite{Sasaki}.

Recall that $d_{1}$ is the distance on $T^{1}\mathbb H^{2}$ defined by
\[
d_{1}(\tilde v,\tilde w)=d\bigl(\tilde v(0),\tilde w(0)\bigr)+d\bigl(\tilde v(1),\tilde w(1)\bigr).
\]

Clearly, both $d_{1}$ and $d_{2}$ are invariant by $G=\mathrm{PSL}(2,\mathbb R)$.

Let $\Gamma$ be a torsion free 
Fuchsian group.
Since $d_{1},d_{2}$ are $G$-invariant, they induce distances, denoted again $d_{1}$ and $d_{2}$,
on $T^{1}S=\Gamma\backslash T^{1}\mathbb H^{2}$ as follows: for $u,v\in T^{1}S$,
\[
\forall i=1,2,\qquad d_{i}(u,v)=\inf_{\gamma\in\Gamma} d_{i}(\tilde u,\gamma\tilde v),
\]
where $\tilde u,\tilde v$ project onto $u,v$.

For $i=1,2$, set
\[
W_{d_{i}}^{ss}(u)=\left\{\,v\in T^{1}S\;\middle|\; \lim\limits_{t\to+\infty} d_{i}(g_{t}u,g_{t}v)=0 \right\},
\]
where $g_{\R}$ is the geodesic flow.

\begin{theorem}\label{equiv_stable_sets}
	\[
	W_{d_{1}}^{ss}(u)=W_{d_{2}}^{ss}(u).
	\]
\end{theorem}

Before proving this theorem, let us prove the following proposition.

\begin{proposition}\label{prop4}
	Let $\tilde u_{0}\in T^{1}\mathbb H^{2}$ such that
	\[
	\tilde u_{0}(0)=i
	\qquad\text{and}\qquad
	\tilde u_{0}(+\infty)=\infty.
	\]
	Let $(\tilde v_{t})_{t\geq 0}$ be a family of vectors in $T^{1}\mathbb H^{2}$.
	Then
	\[
	\lim\limits_{t\to+\infty} d_{1}(\tilde u_{0},\tilde v_{t})=0
	\qquad\Longleftrightarrow\qquad
	\lim\limits_{t\to+\infty} d_{2}(\tilde u_{0},\tilde v_{t})=0.
	\]
\end{proposition}

\begin{proof}

For $t\geq 0$, set
\begin{itemize}
	
	\item$ \alpha_{t}=\alpha_{\tilde u_{0},\tilde v_{t}},
	\qquad
	\beta_{t}=\beta_{\tilde u_{0},\tilde v_{t}},
	$
	\item $\xi_{t}=C_{\tilde u_{0},\tilde v_{t}}(+\infty),
	\qquad
	\xi_{t}'=\tilde v_{t}(+\infty).
	$
\end{itemize}
\medskip

\noindent
$\left( \Longrightarrow \right)$ Suppose $
\lim\limits_{t\to+\infty} d_{1}(\tilde u_{0},\tilde v_{t})=0.
$

Since $
\lim\limits_{t\to+\infty}\tilde v_{t}(0)=i
\text{ and }
\lim\limits_{t\to+\infty}\tilde v_{t}(1)=e\,i,
$
we have $
\lim\limits_{t\to+\infty}\xi_{t}'=\infty.
$

Suppose by contradiction that there exist $\varepsilon>0$ and a sequence $(t_{n})_{n\geq 0}\subset\mathbb R_{+}$
converging to $+\infty$ such that
\begin{equation}\label{eq1}
	\forall n\geq 0,\qquad |\alpha_{t_{n}}-\beta_{t_{n}}|\geq \varepsilon.
\end{equation}

Up to a subsequence, we can suppose $
\lim\limits_{n\to+\infty}\xi_{t_{n}}=\xi.
$

Let $\tilde w\in T^{1}\mathbb H^{2}$ such that
$
\tilde w(0)=i
\text{ and }
\tilde w(+\infty)=\xi,
$

we have
$
\lim\limits_{n\to+\infty} C_{\tilde u_{0},\tilde v_{t_{n}}}=\tilde w,
\lim\limits_{n\to+\infty}g_{s_n}C_{\tilde u_0,\tilde v_{t_n}}=\tilde w, \text{ with } s_n=d(i,\tilde v_{t_n}(0))
\text{ and }
\lim\limits_{n\to+\infty}\tilde v_{t_{n}}=\tilde u_{0}.
$

It follows that
\[
\lim\limits_{n\to+\infty}\alpha_{t_{n}}=\widehat{(\tilde w,\tilde u_{0})}
\qquad\text{and}\qquad
\lim\limits_{n\to+\infty}\beta_{t_{n}}=\widehat{(\tilde w,\tilde u_{0})},
\]
hence
$
\lim\limits_{n\to+\infty}|\alpha_{t_{n}}-\beta_{t_{n}}|=0,
$
which contradicts (\ref{eq1}).

\medskip

\noindent
$\left( \Longleftarrow \right)$ Suppose
$
\lim\limits_{t\to+\infty} d_{2}(\tilde u_{0},\tilde v_{t})=0.
$

Suppose by contradiction that there exist $\varepsilon>0$ and $(t_{n})_{n\geq 0}\subset\mathbb R_{+}$
converging to $+\infty$ such that
\begin{equation}\label{eq2}
	\forall t_{n},\qquad d\bigl(\tilde u_{0}(1),\tilde v_{t_{n}}(1)\bigr)\geq \varepsilon.
\end{equation}

We can suppose
$
\lim\limits_{n\to+\infty}\xi_{t_{n}}=\xi
\text{ and }
\lim\limits_{n\to+\infty}\xi_{t_{n}}'=\xi'.
$

Let $\tilde w\in T^{1}\mathbb H^{2}$ with
$
\tilde w(0)=i
\text{ and }
\tilde w(+\infty)=\xi,
$
and $\tilde w'\in T^{1}\mathbb H^{2}$ with
$
\tilde w'(0)=i
\text{ and } 
\tilde w'(+\infty)=\xi'.
$

Clearly,
$
\lim\limits_{n\to+\infty}\tilde v_{t_{n}}=\tilde w'
\text{ and } 
\lim\limits_{n\to+\infty} C_{\tilde u_{0},\tilde v_{t_{n}}}=\tilde w.
$

It follows that
$
\lim\limits_{n\to+\infty}\alpha_{t_{n}}=\widehat{(\tilde w,\tilde u_{0})}
\text{ and } 
\lim\limits_{n\to+\infty}\beta_{t_{n}}=\widehat{(\tilde w,\tilde w')}.
$ 

Since
$
\lim\limits_{n\to+\infty}|\alpha_{t_{n}}-\beta_{t_{n}}|=0,
$
we obtain
$
\tilde w'=\tilde u_{0}.
$

We deduce that the geodesic segment
$
[\tilde v_{t_{n}}(0),\tilde v_{t_{n}}(1)]
$
converges to the vertical arc
$
[i,e\,i],
$
which contradicts (\ref{eq2}).
\end{proof}

\subsection*{Proof of the Theorem \ref{equiv_stable_sets}}

Let $\tilde u,\tilde v\in T^{1}\mathbb H^{2}$ projecting onto $u,v$ in $T^{1}S$ such that
$v\in W_{d_{i}}^{ss}(u)$.

By definition, there exists $(\gamma_{t})_{t\geq 0}$ in $\Gamma$ such that
\[
\lim\limits_{t\to+\infty} d_{i}\bigl(g_{t}\tilde u,\gamma_{t}g_{t}\tilde v\bigr)=0.
\]

Since $G$ acts on $T^{1}\mathbb H^{2}$ transitively, there exists a sequence of isometries $(f_{t})_{t\geq 0}$ in $G$ such that
\[
g_{t}\tilde u=f_{t}\tilde u_{0}.
\]

Since $G$ acts by isometries on $(T^{1}\mathbb H^{2},d_{i})$, we have
\[
d_{i}\bigl(g_{t}\tilde u,\gamma_{t}g_{t}\tilde v\bigr)
=
d_{i}\bigl(\tilde u_{0},f_{t}^{-1}\gamma_{t}g_{t}\tilde v\bigr).
\]

Applying Proposition \ref{prop4}, we obtain that $v\in W_{d_{j}}^{ss}(u)$, with
$
j\in\{1,2\}.
$
\qed

	
	
	
\begin{thebibliography}{99}
	
	
	\bibitem{Ballmann}
	W. Ballmann, M. Gromov and V. Schroeder, {\it Manifolds of nonpositive curvature}, Progress in Mathematics, 61, Birkh\"auser Boston, Boston, MA, 1985.
	
	\bibitem{Beardon}
	A.~F. Beardon, {\it The geometry of discrete groups}, Graduate Texts in Mathematics, 91, Springer, New York, 1983.
	
	\bibitem{Bellis}
	A.~Bellis,
	\emph{�tude topologique du flot horocyclique: le cas des surfaces g�om�triquement infinies},
	PhD thesis, Institut de Recherche Math�matique de Rennes, 2018.
	
%	\bibitem{BurniolDalBo}
%	S.~Burniol Clotet and F.~Dal'Bo,
%	\emph{The geodesic flow on a hyperbolic surface with a cusp is not expansive},
%	arXiv:2603.24310, 2026.

	\bibitem{Dalbo} 
	F. Dal'Bo, {\it Trajectoires g\'eod\'esiques et horocycliques}, Savoirs Actuels, EDP Sciences, Les Ulis, 2007.

	\bibitem{Sasaki} 
	S. Sasaki, On the differential geometry of tangent bundles of Riemannian manifolds, Tohoku Math. J. (2) {\bf 10} (1958), 338--354.


\end{thebibliography}

\end{document}
